% Section file for Chapter: Twisted Holography and Quantum Gravity
% Result (G1): Perturbative Finiteness
%
% Develops the HS-sewing framework, shadow free energies,
% convergence of the gravitational partition function, the
% Heisenberg prototype, finiteness for the standard landscape,
% and implications for 3d quantum gravity.

\section{Perturbative finiteness}
% label removed: sec:thqg-perturbative-finiteness
\index{perturbative finiteness|textbf}
\index{UV finiteness!twisted holography}

% ----------------------------------------------------------------------
% Local macros: provided only if absent from the ambient manuscript.
% ----------------------------------------------------------------------
\providecommand{\HS}{\operatorname{HS}}
\providecommand{\Tr}{\operatorname{Tr}}
\providecommand{\Det}{\operatorname{Det}}
\providecommand{\Fred}{\operatorname{Fred}}
\providecommand{\rank}{\operatorname{rank}}
\providecommand{\Sym}{\operatorname{Sym}}
\providecommand{\id}{\mathrm{id}}
\providecommand{\ad}{\operatorname{ad}}
\providecommand{\wt}{\operatorname{wt}}
\providecommand{\Aut}{\operatorname{Aut}}
\providecommand{\Res}{\operatorname{Res}}
\providecommand{\Defcyc}{\operatorname{Def}_{\mathrm{cyc}}}
\providecommand{\gAmod}{\mathfrak{g}_{\cA}^{\mathrm{mod}}}
\providecommand{\Gmod}{\mathcal{G}^{\mathrm{mod}}}

The perturbative finiteness of twisted quantum gravity is the foundational analytic result that makes the gravitational shadow programme well defined.  In standard perturbative quantum gravity, the genus-$g$ contribution to the path integral diverges for~$g \geq 2$ (nonrenormalizability of the Einstein--Hilbert action in~$d = 4$).  In the twisted setting, where the boundary chiral algebra~$\cA$ controls the bulk through the holographic modular Koszul datum $\mathcal{H}(T) = (\cA, \cA^!, \mathcal{C}, r(z), \Theta_\cA, \nabla^{\mathrm{hol}})$, the situation is radically different: the Maurer--Cartan equation $D_\cA\Theta_\cA + \tfrac{1}{2}[\Theta_\cA,\Theta_\cA] = 0$ is an \emph{exact} identity, not a truncation of an asymptotic series, and the Fulton--MacPherson compactification provides a geometric regulator that renders every genus-$g$ amplitude finite without the need for UV counterterms.

The section develops this finiteness in six stages: (\S\ref{subsec:thqg-I-hs-sewing-framework}) the HS-sewing framework; (\S\ref{subsec:thqg-I-shadow-free-energies}) shadow free energies and the Bernoulli generating function; (\S\ref{subsec:thqg-I-convergence-grav}) convergence of the gravitational partition function; (\S\ref{subsec:thqg-I-heisenberg-prototype}) the Heisenberg partition function as prototype; (\S\ref{subsec:thqg-I-standard-landscape}) finiteness for the standard landscape; (\S\ref{subsec:thqg-I-3d-gravity}) implications for 3d quantum gravity.

% ======================================================================
%  SUBSECTION 1: THE HS-SEWING FRAMEWORK FOR TWISTED GRAVITY
% ======================================================================

\subsection{The HS-sewing framework for twisted gravity}
% label removed: subsec:thqg-I-hs-sewing-framework
\index{HS-sewing!twisted gravity framework}

The passage from finite-dimensional algebraic data (the OPE coefficients of the boundary chiral algebra~$\cA$) to infinite-genus partition functions requires a bridge between the algebraic and analytic worlds.  This bridge is the HS-sewing condition, which guarantees that the sewing amplitudes on pairs of pants compose to give trace-class operators on the sewing Hilbert space.

\subsubsection{The sewing Hilbert space}

\begin{definition}[Sewing Hilbert space; cf.\ {\cite{costello-renormalization}}]
% label removed: def:thqg-I-sewing-hilbert-space
\index{sewing Hilbert space|textbf}
Let $\cA$ be a positive-energy chiral algebra with conformal grading $\cA = \bigoplus_{n \geq 0} H_n$, where $\dim H_n < \infty$ for all~$n$.  For a sewing parameter $q \in (0,1)$, the \emph{sewing Hilbert space} is the $q$-weighted $\ell^2$-completion
\begin{equation}% label removed: eq:thqg-I-sewing-hilbert
\mathcal{H}_q(\cA) \;:=\; \widehat{\bigoplus}_{n \geq 0} q^{n/2}\,H_n \;=\; \Bigl\{ \sum_{n=0}^\infty v_n \;\Big|\; v_n \in H_n,\; \sum_{n=0}^\infty q^n \|v_n\|^2 < \infty \Bigr\},
\end{equation}
where $\|\cdot\|$ is any fixed Hermitian norm on each finite-dimensional $H_n$.  The inner product is $\langle v, w \rangle_q = \sum_{n \geq 0} q^n \langle v_n, w_n \rangle_{H_n}$.
\end{definition}

\begin{remark}[Independence of norms]
% label removed: rem:thqg-I-norm-independence
Since each $H_n$ is finite-dimensional, different choices of Hermitian norm on $H_n$ give equivalent topologies on $\mathcal{H}_q(\cA)$.  The $q$-weighting is the essential ingredient: it ensures that the completion is sensitive to conformal weight, with high-weight states exponentially suppressed.  For $q \to 1^-$, the Hilbert space degenerates; for $q \to 0^+$, only the vacuum sector $H_0$ survives.
\end{remark}

\begin{lemma}[Completeness and separability; \ClaimStatusProvedHere]
% label removed: lem:thqg-I-completeness
\index{sewing Hilbert space!completeness}
$\mathcal{H}_q(\cA)$ is a separable Hilbert space for every $q \in (0,1)$.  A Cauchy sequence $\{v^{(k)}\}_{k \geq 1}$ converges if and only if $\sum_{n \geq 0} q^n \|v^{(k)}_n - v^{(\ell)}_n\|^2 \to 0$ as $k, \ell \to \infty$.
\end{lemma}

\begin{proof}
Each $H_n$ is finite-dimensional, hence complete.  The $q$-weighted $\ell^2$-completion is a weighted direct sum of Hilbert spaces with weights $q^n$; completeness of the direct sum follows from that of each summand together with the summability condition.  Separability follows from taking a countable dense subset of each $H_n$ (finite-dimensional spaces are separable) and forming finite linear combinations.
\end{proof}

\begin{proposition}[Inclusion maps; \ClaimStatusProvedHere]
% label removed: prop:thqg-I-inclusion-maps
\index{sewing Hilbert space!inclusion}
For $0 < q_1 < q_2 < 1$, the identity map on $\cA$ extends to a bounded inclusion $\iota_{q_1,q_2} \colon \mathcal{H}_{q_2}(\cA) \hookrightarrow \mathcal{H}_{q_1}(\cA)$ with operator norm $\|\iota_{q_1,q_2}\| = 1$.  In particular, for every $v \in \mathcal{H}_{q_2}(\cA)$:
\begin{equation}% label removed: eq:thqg-I-inclusion
\|v\|_{q_1}^2 \;=\; \sum_{n \geq 0} q_1^n \|v_n\|^2 \;\leq\; \sum_{n \geq 0} q_2^n \|v_n\|^2 \;=\; \|v\|_{q_2}^2.
\end{equation}
\end{proposition}

\begin{proof}
Since $q_1 < q_2$, we have $q_1^n \leq q_2^n$ for all $n \geq 0$, giving the inequality termwise.
\end{proof}

\subsubsection{The HS-sewing condition}

\begin{definition}[HS-sewing condition]
% label removed: def:thqg-I-hs-sewing
\index{HS-sewing!definition|textbf}
Let $\cA$ be a positive-energy chiral algebra with OPE structure constants $C^{c,k}_{a,i;\,b,j}$, where $a,b,c$ are conformal weights and $i,j,k$ index bases of $H_a$, $H_b$, $H_c$ respectively.  The \emph{pair-of-pants multiplication operator} at weight $(a,b) \to c$ is the linear map $m^c_{a,b} \colon H_a \otimes H_b \to H_c$ with matrix entries $(m^c_{a,b})_{ij}^k = C^{c,k}_{a,i;\,b,j}$.  The algebra~$\cA$ satisfies the \emph{HS-sewing condition} at parameter $q \in (0,1)$ if
\begin{equation}% label removed: eq:thqg-I-hs-sewing-condition
\boxed{ \sum_{a,b,c \geq 0} q^{a+b+c}\, \|m^c_{a,b}\|_{\HS}^2 \;<\; \infty,}
\end{equation}
where $\|m^c_{a,b}\|_{\HS}^2 = \sum_{i,j,k} |C^{c,k}_{a,i;\,b,j}|^2$ is the Hilbert--Schmidt norm of the multiplication operator.
\end{definition}

\begin{remark}[Physical content]
% label removed: rem:thqg-I-hs-physical
The HS-sewing condition controls the convergence of the sewing construction: a pair-of-pants decomposition of a genus-$g$ surface with $n$ punctures involves $2g - 2 + n$ pairs of pants, and each sewing circle contributes one composition of pair-of-pants operators.  The HS condition guarantees that each such composition is Hilbert--Schmidt, and that the composition of two Hilbert--Schmidt operators is trace class, the key requirement for the trace that computes the partition function.  The $q$-parameter is the sewing modulus $q = e^{-t}$ where $t > 0$ is the collar length.
\end{remark}

\begin{definition}[Hilbert--Schmidt and trace-class norms]
% label removed: def:thqg-I-hs-trace-norms
\index{Hilbert--Schmidt norm}
\index{trace-class!operator}
For a linear map $T \colon V \to W$ between finite-dimensional Hilbert spaces, the Hilbert--Schmidt norm is $\|T\|_{\HS}^2 = \Tr(T^*T) = \sum_{i,j} |T_{ij}|^2$.  An operator $K$ on a Hilbert space $H$ is \emph{trace class} if $\|K\|_1 := \Tr\sqrt{K^*K} < \infty$.  The key composition inequality is:
\begin{equation}% label removed: eq:thqg-I-hs-composition
\|AB\|_1 \;\leq\; \|A\|_{\HS}\,\|B\|_{\HS}
\end{equation}
for Hilbert--Schmidt operators $A, B$.  Hence the composition of two HS operators is trace class, with trace bounded by the product of HS norms.
\end{definition}

\begin{lemma}[Schatten class hierarchy; \ClaimStatusProvedHere]
% label removed: lem:thqg-I-schatten
\index{Schatten class}
Let $H$ be a separable Hilbert space and $S_p(H)$ the $p$-th Schatten class ($p \geq 1$).  Then:
\begin{enumerate}[label=\textup{(\roman*)}]
\item $S_1(H) \subset S_2(H) \subset S_p(H) \subset \mathcal{B}(H)$ for $1 \leq p \leq \infty$, where $\mathcal{B}(H)$ is bounded operators.
\item $S_1(H)$ is a two-sided ideal in $\mathcal{B}(H)$: if $K \in S_1(H)$ and $T \in \mathcal{B}(H)$, then $KT, TK \in S_1(H)$.
\item If $A \in S_p(H)$ and $B \in S_q(H)$ with $1/p + 1/q = 1/r \leq 1$, then $AB \in S_r(H)$.
\end{enumerate}
The sewing application uses $p = q = 2$, $r = 1$: two Hilbert--Schmidt operators compose to trace class.
\end{lemma}

\begin{proof}
Standard functional analysis; see Reed--Simon~\cite{RS78} or Simon~\cite{Simon05}.  Part~(iii) is the generalized H\"older inequality for Schatten classes.
\end{proof}

\begin{proposition}[Pair-of-pants as HS operator; \ClaimStatusProvedElsewhere]
% label removed: prop:thqg-I-pants-hs
\index{pair of pants!HS operator}
Under HS-sewing~\eqref{eq:thqg-I-hs-sewing-condition}, the $q$-weighted multiplication operator $m_q \colon \mathcal{H}_q(\cA) \,\widehat{\otimes}\, \mathcal{H}_q(\cA) \to \mathcal{H}_q(\cA)$ is Hilbert--Schmidt, with
\begin{equation}% label removed: eq:thqg-I-mq-hs-bound
\|m_q\|_{\HS}^2 \;=\; \sum_{a,b,c \geq 0} q^{a+b+c}\,\|m^c_{a,b}\|_{\HS}^2 \;<\; \infty.
\end{equation}
Consequently, every closed amplitude on a genus-$g$ surface obtained by sewing is trace class.
\end{proposition}

\begin{proof}
Direct from the definition: $\|m_q\|_{\HS}^2$ is precisely the left-hand side of~\eqref{eq:thqg-I-hs-sewing-condition}.  A genus-$g$ surface with $n$ punctures is obtained by sewing $2g - 2 + n$ pairs of pants along $3g - 3 + n$ circles.  Each sewing circle contributes one pair-of-pants composition.  By the Hilbert--Schmidt composition inequality~\eqref{eq:thqg-I-hs-composition}, each composition of two HS operators gives a trace-class operator.  Since $g \geq 1$ requires at least two compositions, and trace-class operators form an ideal, the result follows.
\end{proof}

\begin{remark}[Sewing topology]
% label removed: rem:thqg-I-sewing-topology
\index{sewing!topology}
The pair-of-pants decomposition of a genus-$g$ surface $\Sigma_{g,n}$ into $2g - 2 + n$ trinions and $3g - 3 + n$ sewing circles is determined by a pants graph $\Gamma_P$.  Different choices of $\Gamma_P$ give different factorizations of the sewing operator, but the trace (the partition function) is independent of $\Gamma_P$ by the factorization property.  The sewing topology on the space of pants decompositions is captured by the flip graph, whose vertices are pants decompositions and whose edges are elementary moves (Dehn twists and handle slides).  The partition function is invariant under flips by the pentagon and hexagon identities of the chiral algebra.
\end{remark}

\subsubsection{The polynomial-OPE / subexponential-sector criterion}

The fundamental analytic input is the criterion of Theorem~\ref{thm:general-hs-sewing}: polynomial OPE growth and subexponential sector growth together imply HS-sewing for all $q \in (0,1)$.

\begin{theorem}[HS-sewing from polynomial OPE growth; \ClaimStatusProvedElsewhere{} {\textup{Theorem~\ref{thm:general-hs-sewing}}}]
% label removed: thm:thqg-I-hs-sewing-criterion
\index{HS-sewing!polynomial OPE criterion}
Let\/ $\cA$ be a positive-energy chiral algebra satisfying:
\begin{enumerate}[label=\textup{(\alph*)}]
\item \emph{Subexponential sector growth:} $\log\dim H_n = o(n)$ as $n \to \infty$.
\item \emph{Polynomial OPE growth:} There exist constants $K > 0$ and $N \in \mathbb{Z}_{\geq 0}$ such that $|C^{c,k}_{a,i;\,b,j}| \leq K(a + b + c + 1)^N$ for all conformal weights $a,b,c$ and all basis indices $i,j,k$.
\end{enumerate}
Then $\cA$ satisfies the HS-sewing condition for every $q \in (0,1)$:
\begin{equation}% label removed: eq:thqg-I-hs-criterion-bound
\sum_{a,b,c \geq 0} q^{a+b+c}\,\|m^c_{a,b}\|_{\HS}^2 \;\leq\; K^2 \sum_{n \geq 0} P_3(n)\,n^{2N}\, \bigl(\!\dim H_n\bigr)^3\,q^n \;<\; \infty,
\end{equation}
where $P_3(n) = \binom{n+2}{2}$ counts the number of triples $(a,b,c)$ with $a + b + c = n$.
\end{theorem}

\begin{proof}
The HS norm at weight level $(a,b,c)$ satisfies
\[
\|m^c_{a,b}\|_{\HS}^2 \;=\; \sum_{i=1}^{\dim H_a} \sum_{j=1}^{\dim H_b} \sum_{k=1}^{\dim H_c} |C^{c,k}_{a,i;\,b,j}|^2 \;\leq\; \dim H_a \cdot \dim H_b \cdot \dim H_c \cdot K^2(a{+}b{+}c{+}1)^{2N}.
\]
Summing over triples with $a + b + c = n$:
\[
\sum_{\substack{a,b,c \geq 0 \\ a+b+c = n}} q^n\,\|m^c_{a,b}\|_{\HS}^2 \;\leq\; K^2 q^n n^{2N} \sum_{\substack{a,b,c \geq 0 \\ a+b+c = n}} \dim H_a \cdot \dim H_b \cdot \dim H_c.
\]
The inner sum over triples is bounded by $\bigl(\sum_{j=0}^n \dim H_j\bigr)^3$.  By the subexponential growth hypothesis, for every $\varepsilon > 0$ there exists $C_\varepsilon$ with $\dim H_j \leq C_\varepsilon \,e^{\varepsilon j}$.  Hence $\sum_{j=0}^n \dim H_j \leq C_\varepsilon\,n\,e^{\varepsilon n}$.  The total sum is therefore bounded by $K^2 C_\varepsilon^3 \sum_{n \geq 0} n^{2N+3}\, e^{3\varepsilon n}\,q^n$.  Choosing $\varepsilon < |\log q|/4$, the exponential factor $e^{3\varepsilon n} q^n = e^{(3\varepsilon + \log q)n}$ decays geometrically, ensuring convergence of the power series.
\end{proof}

\begin{remark}[Optimality of the bound]
% label removed: rem:thqg-I-hs-optimality
The polynomial-OPE criterion is not optimal: examples with super-polynomial but sub-exponential OPE growth (such as certain completed W-algebras) may still satisfy HS-sewing.  The point is that the criterion covers the entire standard landscape of finitely strongly generated chiral algebras.  The bound~\eqref{eq:thqg-I-hs-criterion-bound} is loose by factors of order $n^3$ due to the crude estimate $\sum_{a+b+c=n} \dim H_a \dim H_b \dim H_c \leq (\sum_{j \leq n} \dim H_j)^3$; tighter bounds are available for specific families.
\end{remark}

\begin{remark}[The gap between HS-sewing and convergence]
% label removed: rem:thqg-I-gap-hs
\index{HS-sewing!necessary vs sufficient}
The HS-sewing condition is \emph{sufficient} for absolute convergence of the genus-$g$ partition function but not \emph{necessary}.  Weaker conditions (e.g., trace-class sewing without the Hilbert--Schmidt factorization) may suffice for finiteness at specific genera.  The HS condition provides a uniform bound across all genera simultaneously, which is the key advantage: a single verification at the level of OPE data yields finiteness at every genus.
\end{remark}

\subsubsection{Explicit HS norms for the standard landscape}

\begin{computation}[HS norms for Heisenberg]
% label removed: comp:thqg-I-hs-heisenberg
\index{Heisenberg!HS norms}
For the Heisenberg algebra $\mathcal{H}_k$ of rank~$1$ at level $k \in \mathbb{C}^{\times}$:
\begin{itemize}
\item $\dim H_n = p(n)$ (partitions of~$n$);
\item OPE: $a(z)a(w) \sim k/(z-w)^2$, giving $C^c_{a,b} = k\,\delta_{a+b,c+1}$ for the fundamental OPE, and all composite OPE grow polynomially with degree $N = 1$ (the Wick contraction is bilinear).
\item Growth: $\log p(n) \sim \pi\sqrt{2n/3}$ (Hardy--Ramanujan), so $\log \dim H_n = O(\sqrt{n})$, which is subexponential.
\end{itemize}
The HS sum factorizes by Wick's theorem into a one-particle contribution:
\begin{equation}% label removed: eq:thqg-I-hs-heisenberg-norm
\|m_q\|_{\HS}^2 \;=\; \prod_{n=1}^\infty \frac{1}{(1-q^n)^{2|k|^2}} \;=\; \eta(q)^{-2|k|^2},
\end{equation}
which is finite for all $q \in (0,1)$.  The one-particle operator has eigenvalues $q^n$, giving $\|T_q\|_1 = \sum_{n=1}^\infty q^n = q/(1-q)$, which is trace class for all $q < 1$.

The factorization into the Dedekind eta function is the consequence of the boson-fermion correspondence: the multi-particle Fock space partition function is the exponential of the one-particle trace, and the Dedekind eta function $\eta(q) = q^{1/24}\prod_{n=1}^\infty (1-q^n)$ is the generating function for partitions.  Explicitly, the second quantization identity gives
\begin{equation}% label removed: eq:thqg-I-hs-heisenberg-factorization
\sum_{n=0}^\infty p(n)^2 q^n \;=\; \prod_{m=1}^\infty \frac{1}{(1-q^m)^2},
\end{equation}
and the HS norm is $|k|^2$ copies of this product.
\end{computation}

\begin{computation}[HS norms for affine Kac--Moody]
% label removed: comp:thqg-I-hs-affine
\index{affine Kac--Moody!HS norms}
For $\widehat{\fg}_k$ at non-critical level $k \neq -h^\vee$, with $\fg$ simple of dimension~$d$ and dual Coxeter number $h^\vee$:
\begin{itemize}
\item $\dim H_n \sim C\,n^{(d-\rank)/2}\,e^{\pi\sqrt{2dn/3}}$ (Kac--Peterson), so $\log \dim H_n = O(\sqrt{n})$, subexponential.
\item OPE: $J^a(z)J^b(w) \sim k\delta^{ab}/(z-w)^2 + f^{ab}{}_c J^c(w)/(z-w)$, polynomial of degree $N = 1$.
\item The HS norm satisfies
  \begin{equation}% label removed: eq:thqg-I-hs-affine-norm
  \|m_q\|_{\HS}^2 \;\leq\; C_k^2 \prod_{n=1}^\infty \frac{1}{(1-q^n)^{2d}} \;=\; C_k^2\,\eta(q)^{-2d}
  \end{equation}
  where $C_k$ depends on $k$ and the structure constants of $\fg$.
\end{itemize}
The key estimate for $C_k$ is obtained from the OPE normalization.  The structure constants of $\widehat{\fg}_k$ at weight level $n$ grow as $|C^c_{a,b}| \leq \|f\|_\infty \cdot (k+h^\vee)$, where $\|f\|_\infty = \max_{a,b,c} |f^{ab}{}_c|$ is the maximum structure constant.  For $\fg = \mathfrak{sl}_N$, one has $\|f\|_\infty = 1$ in the Chevalley basis, giving $C_k = |k + h^\vee|$.
\end{computation}

\begin{computation}[HS norms for Virasoro]
% label removed: comp:thqg-I-hs-virasoro
\index{Virasoro!HS norms}
For the Virasoro algebra $\mathrm{Vir}_c$ at central charge $c \in \mathbb{C}$:
\begin{itemize}
\item $\dim H_n = p(n)$ (Verma module dimensions);
\item OPE: $T(z)T(w) \sim c/(2(z-w)^4) + 2T(w)/(z-w)^2 + \partial T(w)/(z-w)$, polynomial of degree $N = 2$ (the $(z-w)^{-4}$ pole gives OPE coefficients growing at most quadratically in conformal weight).
\item Growth: $\log p(n) \sim \pi\sqrt{2n/3}$, subexponential.
\item The HS sum converges:
  \begin{equation}% label removed: eq:thqg-I-hs-virasoro-norm
  \|m_q\|_{\HS}^2 \;\leq\; C_c^2 \sum_{n \geq 0} p(n)^3\,n^{4}\,q^n \;<\; \infty \quad\text{for all } q \in (0,1).
  \end{equation}
\end{itemize}
The quartic growth of $n^4$ in the bound reflects the quartic OPE pole of $T(z)T(w)$.  For the Virasoro algebra, the structure constants $C^c_{a,b}$ at weight level $n = a + b - c$ satisfy $|C^c_{a,b}| \leq P_2(n) := (n+1)(n+2)/2$ when $c$ is bounded; the precise bound involves the central charge through $|c|/2$ from the fourth-order pole.
\end{computation}

\begin{computation}[HS norms for $\beta\gamma$ system]
% label removed: comp:thqg-I-hs-betagamma
\index{beta-gamma system@$\beta\gamma$ system!HS norms}
For the $\beta\gamma$ system of conformal weights $(1,0)$:
\begin{itemize}
\item $\dim H_n = p(n)^2$ (two independent partition functions, one for $\beta$ modes and one for $\gamma$ modes).
\item OPE: $\beta(z)\gamma(w) \sim 1/(z-w)$, a simple pole.  All composite OPE grow polynomially with degree $N = 1$.
\item Growth: $\log p(n)^2 = 2\log p(n) \sim 2\pi\sqrt{2n/3}$, subexponential.
\item The HS norm:
  \begin{equation}% label removed: eq:thqg-I-hs-betagamma-norm
  \|m_q\|_{\HS}^2 \;\leq\; C^2\,\eta(q)^{-4} \;<\; \infty \quad\text{for all } q \in (0,1).
  \end{equation}
\end{itemize}
The exponent $-4$ in $\eta(q)^{-4}$ reflects the two bosonic degrees of freedom ($\beta$ and $\gamma$), each contributing $\eta(q)^{-2}$.
\end{computation}

\begin{computation}[HS norms for $\mathcal{W}_N$]
% label removed: comp:thqg-I-hs-wn
\index{W-algebra@$\mathcal{W}$-algebra!HS norms}
For the principal $\mathcal{W}$-algebra $\mathcal{W}_k(\mathfrak{sl}_N)$ at generic level:
\begin{itemize}
\item $\dim H_n$ is the number of $N$-colored partitions of~$n$ with colors of weight $2, 3, \ldots, N$; the growth is $\log \dim H_n = O(n^{(N-1)/N})$, subexponential for all~$N$.
\item OPE growth: the maximal singular order of the $W^{(s)}$ self-OPE is $2s$, so $N_{\max} = 2N$ for the $W^{(N)}$ field.  All OPE coefficients grow at most polynomially with degree $\leq 2N$.
\item Convergence:
  \begin{equation}% label removed: eq:thqg-I-hs-wn-norm
  \|m_q\|_{\HS}^2 \;\leq\; K_N^2(c) \sum_{n \geq 0} (\dim H_n)^3\,n^{4N}\,q^n \;<\; \infty \quad\text{for all } q \in (0,1).
  \end{equation}
\end{itemize}
The $N$-colored partition function has generating series $\prod_{s=2}^N \prod_{m=s}^\infty (1-q^m)^{-1}$, and its growth exponent is controlled by the Meinardus theorem: $\log(\dim H_n) \sim C_N \cdot n^{(N-1)/N}$ with $C_N = N \cdot (\Gamma(1/N) \zeta(1+1/N))^{N/(N+1)} / (N+1)$.  For $N = 2$ (Virasoro), $C_2 = \pi\sqrt{2/3}$; for $N = 3$ ($\mathcal{W}_3$), the growth is $O(n^{2/3})$.
\end{computation}

\begin{computation}[HS norms for lattice VOAs]
% label removed: comp:thqg-I-hs-lattice
\index{lattice VOA!HS norms}
For the lattice vertex algebra $V_\Lambda$ of an even integral lattice $\Lambda$ of rank~$r$:
\begin{itemize}
\item $\dim H_n = |\Lambda^*/\Lambda| \cdot p_r(n)$, where $p_r(n)$ is the $r$-colored partition function.  Growth: $\log \dim H_n = O(\sqrt{n})$, subexponential.
\item OPE: inherited from rank-$r$ Heisenberg plus lattice vertex operators $V_\alpha(z)V_\beta(w) \sim (z-w)^{\langle\alpha,\beta\rangle}:\!V_{\alpha+\beta}(w)\!:$; the OPE growth is polynomial of degree $N = \max(r, \max_{\alpha \in \Lambda}|\langle\alpha,\alpha\rangle|/2)$.
\item The HS norm factors as
  \begin{equation}% label removed: eq:thqg-I-hs-lattice-norm
  \|m_q\|_{\HS}^2 \;\leq\; C_\Lambda^2\,\eta(q)^{-2r} \,\Theta_\Lambda(q)^2 \;<\; \infty \quad\text{for all } q \in (0,1),
  \end{equation}
  where $\Theta_\Lambda(q) = \sum_{\alpha \in \Lambda} q^{\langle\alpha,\alpha\rangle/2}$ is the lattice theta function.
\end{itemize}
The factorization reflects the tensor-product structure $V_\Lambda = \mathcal{H}^{\otimes r} \otimes \mathbb{C}[\Lambda]$: the Heisenberg part contributes $\eta(q)^{-2r}$ and the lattice group ring contributes $\Theta_\Lambda(q)^2$ (one factor from the bra side, one from the ket side of the HS norm).
\end{computation}

\begin{computation}[HS norms for free fermions]
% label removed: comp:thqg-I-hs-fermion
\index{free fermions!HS norms}
For the free-fermion VOA of rank~$n$ (i.e., $n$ chiral fermions $\psi_1, \ldots, \psi_n$ of conformal weight $1/2$):
\begin{itemize}
\item $\dim H_m = \binom{n}{m}$ for $m \leq n$, $\dim H_m = 0$ for $m > n$ when counting weight-$m/2$ states; more precisely, the full Fock space has $\dim H_m$ growing as partitions into distinct parts, with generating function $\prod_{j=1}^n (1 + q^{j-1/2})$.
\item OPE: $\psi_i(z)\psi_j(w) \sim \delta_{ij}/(z-w)$, simple pole.  OPE degree $N = 0$ (the structure constants are $0$ or $\pm 1$).
\item The HS norm:
  \begin{equation}% label removed: eq:thqg-I-hs-fermion-norm
  \|m_q\|_{\HS}^2 \;\leq\; \prod_{j=1}^{n} (1 + q^{j-1/2})^2 \;<\; \infty \quad\text{for all } q \in (0,1).
  \end{equation}
  This bound is finite because it is a finite product of bounded terms.
\end{itemize}
\end{computation}

\begin{table}[ht]
\centering
\caption{HS-sewing growth rates for the standard landscape}
% label removed: tab:thqg-I-hs-growth
\index{HS-sewing!growth rates}
\renewcommand{\arraystretch}{1.4}
{\footnotesize
\begin{tabular}{|l|c|c|c|c|}
\hline
\textbf{Family $\cA$} & $\boldsymbol{\dim H_n}$ & \textbf{OPE growth $N$} & \textbf{Sector growth} & \textbf{HS bound} \\
\hline
$\mathcal{H}_k$ (rank 1) & $p(n)$ & $1$ & $O(\sqrt{n})$ & $\eta(q)^{-2|k|^2}$ \\
\hline
$\widehat{\fg}_k$ & $\sim e^{\pi\sqrt{2dn/3}}$ & $1$ & $O(\sqrt{n})$ & $C_k^2\,\eta(q)^{-2d}$ \\
\hline
$\mathrm{Vir}_c$ & $p(n)$ & $2$ & $O(\sqrt{n})$ & $C_c^2\sum p(n)^3 n^4 q^n$ \\
\hline
$\beta\gamma$ system & $p(n)^2$ & $1$ & $O(\sqrt{n})$ & $\eta(q)^{-4}$ \\
\hline
$\mathcal{W}_N$ & $N$-colored parts & $2N$ & $O(n^{(N-1)/N})$ & $K_N^2(c)\sum (\dim H_n)^3 n^{4N} q^n$ \\
\hline
$V_\Lambda$ (rank $r$) & $|\Lambda^*/\Lambda|\,p_r(n)$ & $r$ & $O(\sqrt{n})$ & $C_\Lambda^2\,\eta(q)^{-2r}\Theta_\Lambda^2$ \\
\hline
Free fermions & $\leq 2^n$ & $0$ & polynomial & $\prod (1+q^{j-1/2})^2$ \\
\hline
\end{tabular}}
\end{table}

\subsubsection{Connection to trace-class operators on sewing Hilbert spaces}

\begin{proposition}[Trace-class amplitudes; \ClaimStatusProvedElsewhere]
% label removed: prop:thqg-I-trace-class-amplitudes
\index{trace class!sewing amplitudes}
Let\/ $\cA$ satisfy the HS-sewing condition at parameter $q$.  Let $\Sigma_{g,n}$ be a genus-$g$ surface with $n$ punctures and a pants decomposition $\mathcal{P}$ with collar parameters $q_1, \ldots, q_{3g-3+n} \in (0,q]$.  Then the sewing amplitude
\begin{equation}% label removed: eq:thqg-I-sewing-amplitude
Z_{\mathcal{P}}(\cA;\, q_1, \ldots, q_{3g-3+n}) \;=\; \Tr_{\mathcal{H}_q(\cA)^{\otimes n}} \Bigl( \prod_{\text{sewing circles}} m_{q_i} \Bigr)
\end{equation}
is well defined and absolutely convergent.  The amplitude depends on the sewing moduli $(q_1, \ldots, q_{3g-3+n})$ but not on the choice of pants decomposition (by the factorization property of the chiral algebra).
\end{proposition}

\begin{proof}
Each pants composition $m_{q_i}$ is Hilbert--Schmidt by Proposition~\ref{prop:thqg-I-pants-hs}.  A genus-$g$ surface requires $2g - 2 + n$ pairs of pants and $3g - 3 + n$ sewing circles.  The sewing amplitude is a trace of a composition of $3g - 3 + n$ Hilbert--Schmidt operators.  Since $3g - 3 + n \geq 2$ for $g \geq 1$ (and $\geq 3$ for $n = 0$), at least one pair of consecutive HS operators composes to give a trace-class operator.  The remaining HS operators are bounded, and trace class times bounded is still trace class (Lemma~\ref{lem:thqg-I-schatten}(ii)).  Hence the trace converges absolutely.  Independence from the pants decomposition follows from the associativity and factorization axioms of the chiral algebra~\cite{BD04, FBZ04}.
\end{proof}

\begin{proposition}[Uniform bound on the trace; \ClaimStatusProvedHere]
% label removed: prop:thqg-I-uniform-trace-bound
\index{trace class!uniform bound}
Under HS-sewing at parameter $q$, the sewing amplitude on a genus-$g$ surface with $n$ punctures satisfies the uniform bound
\begin{equation}% label removed: eq:thqg-I-uniform-trace
|Z_{\mathcal{P}}(\cA;\,q_1,\ldots,q_{3g-3+n})| \;\leq\; \|m_q\|_{\HS}^{3g-3+n}
\end{equation}
for all $q_1, \ldots, q_{3g-3+n} \in (0,q]$.  In particular, the bound is uniform in the sewing moduli.
\end{proposition}

\begin{proof}
Each factor $m_{q_i}$ satisfies $\|m_{q_i}\|_{\HS} \leq \|m_q\|_{\HS}$ (since $q_i \leq q$ and the HS norm is monotone increasing in $q$).  The trace of a composition of operators satisfies $|\Tr(A_1 \cdots A_k)| \leq \|A_1\|_{\HS}\cdots\|A_k\|_{\HS}$ when $k \geq 2$ (by iterated application of the HS composition inequality and the trace bound $|\Tr(K)| \leq \|K\|_1$).
\end{proof}

\begin{corollary}[Genus-$g$ partition function]
% label removed: cor:thqg-I-genus-g-partition
\index{partition function!genus $g$}
For $\cA$ satisfying HS-sewing, the genus-$g$ partition function
\begin{equation}% label removed: eq:thqg-I-Z-g
Z_g(\cA) \;:=\; \int_{\overline{\mathcal{M}}_g} Z_{\mathcal{P}}(\cA;\,\mathbf{q}) \;\mu_{\mathrm{WP}}
\end{equation}
is well defined as an integral over the Deligne--Mumford compactification, where $\mu_{\mathrm{WP}}$ denotes the Weil--Petersson volume form.
\end{corollary}

\begin{proof}
The trace-class norm provides a pointwise bound on the integrand: $|Z_{\mathcal{P}}(\cA;\,\mathbf{q})| \leq \|T_{\mathbf{q}}\|_1$ where $T_{\mathbf{q}}$ is the sewing operator.  By the HS-bound (Proposition~\ref{prop:thqg-I-uniform-trace-bound}), this norm is bounded uniformly in the interior of the moduli space.  Near the boundary, the Fulton--MacPherson compactification provides the regulator: the sewing moduli $q_i$ vanish on the boundary divisors, and the $q$-weighted HS norms degenerate to zero.  The integral over the compact orbifold $\overline{\mathcal{M}}_g$ is therefore finite.
\end{proof}

\begin{remark}[The Deligne--Mumford boundary]
% label removed: rem:thqg-I-dm-boundary
\index{Deligne--Mumford!boundary}
The boundary $\partial\overline{\mathcal{M}}_g = \overline{\mathcal{M}}_g \setminus \mathcal{M}_g$ consists of stable curves with nodes.  Near a boundary stratum $\delta_\Gamma$ indexed by a dual graph $\Gamma$, the sewing moduli satisfy $q_e \to 0$ for each edge $e$ of $\Gamma$.  The integrand $Z_{\mathcal{P}}(\cA;\,\mathbf{q})$ decomposes as a product of amplitudes on the irreducible components of the stable curve, multiplied by propagator insertions at the nodes.  The vanishing of $q_e$ at the boundary ensures that the propagator insertions are trace class with trace-class norm tending to zero, giving the regularity needed for integrability over $\overline{\mathcal{M}}_g$.
\end{remark}


% ======================================================================
%  SUBSECTION 2: SHADOW FREE ENERGIES AND THE BERNOULLI GENERATING FUNCTION
% ======================================================================

\subsection{Shadow free energies and the Bernoulli generating function}
% label removed: subsec:thqg-I-shadow-free-energies
\index{shadow free energy|textbf}
\index{Bernoulli numbers!generating function}

The scalar projection of the Maurer--Cartan element $\Theta_\cA$ onto genus~$g$ is the shadow free energy $F_g(\cA)$.  This is the arity-$0$, genus-$g$ component of the shadow algebra $\cA^{\mathrm{sh}} = H_\bullet(\Defcyc^{\mathrm{mod}}(\cA))$ (Definition~\ref{def:shadow-algebra}).

\subsubsection{Definition of the shadow free energy}

\begin{definition}[Shadow free energy]
% label removed: def:thqg-I-shadow-free-energy
\index{shadow free energy!definition|textbf}
For a modular Koszul chiral algebra $\cA$ with universal Maurer--Cartan class $\Theta_\cA \in \mathrm{MC}(\gAmod)$ (Theorem~\ref*{V1-thm:mc2-bar-intrinsic}), the \emph{genus-$g$ shadow free energy} is the degree-zero shadow representation:
\begin{equation}% label removed: eq:thqg-I-Fg-def
F_g(\cA) \;:=\; \int_{\overline{\mathcal{M}}_g} \operatorname{Sh}_{g,0}(\Theta_\cA) \;=\; \int_{\overline{\mathcal{M}}_g} \kappa(\cA) \cdot \lambda_g,
\end{equation}
where the second equality uses the universality theorem (Theorem~\ref{thm:modular-characteristic}(i)): $\mathrm{obs}_g(\cA) = \kappa(\cA) \cdot \lambda_g$ in $H^{2g}(\overline{\mathcal{M}}_g)$, and $\lambda_g$ is the Hodge class $c_g(\mathbb{E})$.
\end{definition}

\begin{remark}[Relation to the shadow algebra]
% label removed: rem:thqg-I-shadow-algebra-relation
The shadow free energy $F_g(\cA)$ is the $(r,g) = (0,g)$ component of the shadow algebra $\cA^{\mathrm{sh}} = H_\bullet(\Defcyc^{\mathrm{mod}}(\cA))$ in the bigrading by arity~$r$ and genus~$g$.  The higher-arity shadows at genus $g$ (the cubic shadow $\mathfrak{C}_g$ at $r = 3$, the quartic resonance class $\mathfrak{Q}_g$ at $r = 4$) encode the gravitational interactions beyond the free-energy level.  The shadow free energy is the simplest invariant, controlled entirely by the modular characteristic $\kappa(\cA)$.
\end{remark}

\begin{remark}[Physical interpretation: vacuum energy]
% label removed: rem:thqg-I-vacuum-energy
\index{shadow free energy!physical interpretation}
The shadow free energy $F_g(\cA)$ is the genus-$g$ vacuum energy of the twisted holographic theory: it is the contribution to the path integral from genus-$g$ worldsheets with no external insertions.  In string theory language, $F_g$ is the genus-$g$ cosmological constant.  The crucial difference from ordinary string theory is that $F_g$ is determined \emph{exactly} by the modular characteristic $\kappa(\cA)$ through the Faber--Pandharipande integral: there is no moduli-dependent correction, and no need for renormalization.
\end{remark}

\subsubsection{The Faber--Pandharipande integral}

The key input from algebraic geometry is the evaluation of the Hodge integral:

\begin{proposition}[Faber--Pandharipande coefficient; \ClaimStatusProvedElsewhere]
% label removed: prop:thqg-I-faber-pandharipande
\index{Faber--Pandharipande!coefficient}
\index{Hodge integral}
The integral of the top Hodge class over the moduli space of genus-$g$ curves is
\begin{equation}% label removed: eq:thqg-I-fp-integral
\lambda_g^{\mathrm{FP}} \;:=\; \int_{\overline{\mathcal{M}}_g} \lambda_g \;=\; \frac{2^{2g-1} - 1}{2^{2g-1}} \cdot \frac{|B_{2g}|}{(2g)!}\,,
\end{equation}
where $B_{2g}$ is the $2g$-th Bernoulli number.  This formula is due to Faber--Pandharipande~\cite{FP00}.
\end{proposition}

\begin{proof}[Proof sketch]
The integral $\int_{\overline{\mathcal{M}}_g} \lambda_g$ is computed by localization on $\overline{\mathcal{M}}_g$ using the Mumford relation $\lambda_g = (-1)^g \kappa_{2g-2}/(2g-2)!$ and the Witten--Kontsevich theorem for intersection numbers of $\psi$-classes.  The Bernoulli numbers arise through the Euler--Maclaurin formula applied to the localization integral.  A detailed proof is in Faber--Pandharipande~\cite{FP00}; the formula also follows from the string equation and the topological recursion relation.
\end{proof}

\begin{remark}[Bernoulli numbers]
% label removed: rem:thqg-I-bernoulli
\index{Bernoulli numbers!values}
The first ten Bernoulli numbers: $B_0 = 1$, $B_1 = -1/2$, $B_2 = 1/6$, $B_4 = -1/30$, $B_6 = 1/42$, $B_8 = -1/30$, $B_{10} = 5/66$, $B_{12} = -691/2730$, $B_{14} = 7/6$, $B_{16} = -3617/510$, $B_{18} = 43867/798$, $B_{20} = -174611/330$.  The absolute values grow super-exponentially: $|B_{2g}| \sim 2 \cdot (2g)! / (2\pi)^{2g}$ (Euler's formula).  The alternating signs $(-1)^{g+1} B_{2g} > 0$ reflect the alternation of Hodge classes.
\end{remark}

\begin{proposition}[Shadow free energy in closed form; \ClaimStatusProvedElsewhere{} {\textup{Theorem~\ref{thm:modular-characteristic}}}]
% label removed: prop:thqg-I-Fg-closed-form
\index{shadow free energy!closed form}
For any modular Koszul chiral algebra $\cA$:
\begin{equation}% label removed: eq:thqg-I-Fg-formula
\boxed{ F_g(\cA) \;=\; \kappa(\cA) \cdot \frac{2^{2g-1}-1}{2^{2g-1}} \cdot \frac{|B_{2g}|}{(2g)!}\,. }
\end{equation}
In particular, $F_g$ depends only on $g$ and the single scalar $\kappa(\cA)$.
\end{proposition}

\begin{proof}
This is the combination of genus universality ($\mathrm{obs}_g = \kappa \cdot \lambda_g$, cf.\ Theorem~\ref{thm:genus-universality}) with the Faber--Pandharipande evaluation of the Hodge integral (Proposition~\ref{prop:thqg-I-faber-pandharipande}).
\end{proof}

\begin{remark}[Universality]
% label removed: rem:thqg-I-universality
\index{shadow free energy!universality}
The universality of the formula~\eqref{eq:thqg-I-Fg-formula} is striking: the genus-$g$ free energy of any modular Koszul chiral algebra depends only on the single scalar $\kappa(\cA)$, regardless of the detailed OPE structure.  The Heisenberg at level $k$, the affine $\widehat{\mathfrak{sl}}_2$ at level $k$, and the Virasoro at central charge $c$ all have the same functional form $F_g = \kappa \cdot \lambda_g^{\mathrm{FP}}$, differing only in the value of $\kappa$.  This is the shadow-level manifestation of the universality of the modular characteristic (Theorem~D): all nonlinear corrections are encoded in the higher-arity shadows, not in the scalar free energy.
\end{remark}

\subsubsection{The $\hat{A}$-genus generating function}

\begin{theorem}[Generating function for shadow free energies; \ClaimStatusProvedElsewhere{} {\textup{Theorem~\ref{thm:modular-characteristic}(ii)}}]
% label removed: thm:thqg-I-generating-function
\index{A-hat genus@$\hat{A}$-genus!generating function|textbf}
The formal generating series of shadow free energies has the closed-form expression:
\begin{equation}% label removed: eq:thqg-I-gf-ahat
\boxed{ \sum_{g=1}^{\infty} F_g(\cA)\,x^{2g} \;=\; \kappa(\cA) \cdot \left( \frac{x/2}{\sin(x/2)} - 1 \right). }
\end{equation}
Equivalently, with the convention $F_0(\cA) = \kappa(\cA)$:
\begin{equation}% label removed: eq:thqg-I-gf-ahat-alt
\sum_{g=0}^{\infty} F_g(\cA)\,x^{2g} \;=\; \kappa(\cA) \cdot \frac{x/2}{\sin(x/2)}.
\end{equation}
\end{theorem}

\begin{proof}
We use the standard generating function for Bernoulli numbers.  Recall that $\frac{t}{e^t - 1} = \sum_{n=0}^{\infty} \frac{B_n}{n!}\,t^n$.  The shadow free energy uses $\lambda_g^{\mathrm{FP}} = \frac{2^{2g-1}-1}{2^{2g-1}}\cdot\frac{|B_{2g}|}{(2g)!}$ with absolute values of Bernoulli numbers.  The generating function is therefore
\begin{align}
\sum_{g=1}^{\infty} F_g\,x^{2g} &\;=\; \kappa \sum_{g=1}^{\infty} \frac{2^{2g-1}-1}{2^{2g-1}} \cdot \frac{|B_{2g}|}{(2g)!}\,x^{2g} \notag\\
&\;=\; \kappa \left( \frac{x/2}{\sin(x/2)} - 1 \right), % label removed: eq:thqg-I-gf-derivation
\end{align}
where the second equality uses the well-known identity
\begin{equation}% label removed: eq:thqg-I-cosecant-bernoulli
\frac{x/2}{\sin(x/2)} \;=\; 1 + \sum_{g=1}^{\infty} \frac{(2^{2g-1}-1)|B_{2g}|}{2^{2g-1}\cdot(2g)!}\,x^{2g},
\end{equation}
obtainable from the Bernoulli generating function by the substitution $t = ix$ and the relation $\sin(x/2) = (e^{ix/2} - e^{-ix/2})/(2i)$.

The derivation proceeds as follows.  The Bernoulli generating function gives $\frac{t}{e^t - 1} = 1 - t/2 + \sum_{g=1}^\infty \frac{B_{2g}}{(2g)!} t^{2g}$.  Setting $t = ix$:
\[
\frac{ix}{e^{ix} - 1} \;=\; 1 - \frac{ix}{2} + \sum_{g=1}^\infty \frac{B_{2g}}{(2g)!}(ix)^{2g} \;=\; 1 - \frac{ix}{2} + \sum_{g=1}^\infty \frac{(-1)^g B_{2g}}{(2g)!} x^{2g}.
\]
Since $(-1)^g B_{2g} = -|B_{2g}|$ for $g \geq 1$, we have
\[
\frac{ix}{e^{ix}-1} \;=\; 1 - \frac{ix}{2} - \sum_{g=1}^\infty \frac{|B_{2g}|}{(2g)!}x^{2g}.
\]
Combining with the identity $\frac{ix}{e^{ix}-1} + \frac{ix}{2} = \frac{ix/2}{\tanh(ix/2)} = \frac{x/2}{\tan(x/2)}$, and using $\frac{x/2}{\sin(x/2)} = \frac{x/2}{\tan(x/2)} + \frac{x/2\,(1-\cos(x/2))}{\sin(x/2)\cos(x/2)}$, we arrive at~\eqref{eq:thqg-I-cosecant-bernoulli} by the standard manipulation of the partial fraction expansion $\csc(z) = 1/z + 2z\sum_{n=1}^\infty (-1)^n/(z^2 - n^2\pi^2)$.

More directly: the function $\frac{x/2}{\sin(x/2)}$ is even and analytic at $x = 0$ with value $1$.  Its Taylor expansion $\sum_{g=0}^\infty a_{2g} x^{2g}$ has coefficients $a_{2g} = \frac{(2^{2g-1}-1)|B_{2g}|}{2^{2g-1}(2g)!}$, which is verified by comparing with the standard expansion of $x/\sin(x) = \sum_{g=0}^\infty \frac{(2-2^{2-2g})|B_{2g}|}{(2g)!} x^{2g}$ and substituting $x \mapsto x/2$.
\end{proof}

\begin{remark}[Two generating functions]
% label removed: rem:thqg-I-two-gfs
There are two natural generating functions in play.  The \emph{algebraic} generating function $\sum_g F_g\,x^{2g} = \kappa(x/(2\sin(x/2)) - 1)$ uses the real variable $x$ and converges for $|x| < 2\pi$.  The \emph{physical} generating function $\sum_g F_g\,\hbar^g$ uses the genus-counting parameter~$\hbar$ and converges for $|\hbar| < 4\pi^2$.  The relation between the two is $\hbar = x^2$: the algebraic variable $x$ is the ``square root'' of the genus-counting parameter.
\end{remark}

\begin{remark}[Connection to the $\hat{A}$-genus]
% label removed: rem:thqg-I-ahat-connection
\index{A-hat genus@$\hat{A}$-genus!connection to shadow free energy}
The generating function $x/(2\sin(x/2))$ is the square root of the Todd genus $\mathrm{td}(x) = x/(1-e^{-x})$ and is closely related to the $\hat{A}$-genus $\hat{A}(x) = (x/2)/\sinh(x/2)$.  The relation is $x/(2\sin(x/2)) = \hat{A}(ix)$, a Wick rotation.  This Wick rotation corresponds to the passage from the Lorentzian $\hat{A}$-genus (which controls index theory on spin manifolds) to the Euclidean generating function (which controls genus integrals on moduli spaces).  The shadow free energy thus has a natural interpretation as a ``Wick-rotated $\hat{A}$-genus'' of the modular characteristic.
\end{remark}

\subsubsection{Absolute convergence from Bernoulli asymptotics}

\begin{theorem}[Absolute convergence of the generating function; \ClaimStatusProvedHere]
% label removed: thm:thqg-I-absolute-convergence
\index{absolute convergence!shadow free energy}
The genus-counting series $\sum_{g=1}^{\infty} |F_g(\cA)|\,|\hbar|^g$ converges absolutely for $|\hbar| < 4\pi^2$.  The generating function $\sum_{g=1}^{\infty} F_g(\cA)\,\hbar^g$ extends to a meromorphic function on $\mathbb{C}$ via the closed form $\kappa \cdot (\frac{\sqrt{\hbar}/2}{\sin(\sqrt{\hbar}/2)} - 1)$.
\end{theorem}

\begin{proof}
The asymptotic behavior of Bernoulli numbers is controlled by Euler's formula $|B_{2g}| = 2 \cdot \zeta(2g) \cdot (2g)! / (2\pi)^{2g}$, where $\zeta(2g) \to 1$ as $g \to \infty$.  Therefore
\begin{align}
|F_g(\cA)| &\;=\; |\kappa(\cA)| \cdot \frac{2^{2g-1}-1}{2^{2g-1}} \cdot \frac{|B_{2g}|}{(2g)!} \notag \\
&\;=\; |\kappa(\cA)| \cdot \frac{2^{2g-1}-1}{2^{2g-1}} \cdot \frac{2\,\zeta(2g)}{(2\pi)^{2g}} \notag \\
&\;\sim\; \frac{2|\kappa(\cA)|}{(2\pi)^{2g}}. % label removed: eq:thqg-I-Fg-asymptotics
\end{align}
More precisely, for $g \geq 2$:
\[
\frac{2|\kappa|}{(2\pi)^{2g}} \;\leq\; |F_g| \;\leq\; \frac{2|\kappa|}{(2\pi)^{2g}} \cdot \frac{1}{1 - 2^{1-2g}}.
\]
The ratio test gives $|F_{g+1}/F_g| \to 1/(2\pi)^2$ as $g \to \infty$, so the radius of convergence of $\sum F_g\,\hbar^g$ is $(2\pi)^2 = 4\pi^2$.  The closed form $\kappa \cdot (\frac{\sqrt{\hbar}/2}{\sin(\sqrt{\hbar}/2)} - 1)$ is meromorphic on $\mathbb{C}$ with poles at $\hbar = 4\pi^2 n^2$ ($n \geq 1$).

The pole structure is determined by the zeros of $\sin(\sqrt{\hbar}/2)$, which occur at $\sqrt{\hbar}/2 = n\pi$, i.e., $\hbar = 4\pi^2 n^2$.  Near $\hbar = 4\pi^2 n^2$, the function has a simple pole with residue:
\[
\Res_{\hbar = 4\pi^2 n^2} \kappa \cdot \frac{\sqrt{\hbar}/2}{\sin(\sqrt{\hbar}/2)} \;=\; \kappa \cdot \frac{n\pi}{\cos(n\pi)} \cdot \frac{1}{2\sqrt{4\pi^2 n^2}} \cdot 2\sqrt{4\pi^2 n^2} \;=\; \kappa \cdot (-1)^{n+1} \cdot 2n\pi.
\]
\end{proof}

\begin{remark}[Factorially divergent vs.\ convergent]
% label removed: rem:thqg-I-factorial-convergent
\index{factorial growth!absence of}
In standard perturbative quantum field theory, the genus-$g$ contribution typically grows as $(2g)!$, reflecting the factorial number of Feynman diagrams.  The shadow free energy $F_g \sim 2\kappa/(2\pi)^{2g}$ decays exponentially in $g$, not factorially.  The series $\sum F_g \hbar^g$ is convergent, not merely asymptotic.  The reason is that the shadow free energy is an integral over the compact moduli space $\overline{\mathcal{M}}_g$ of a smooth cohomology class, not a sum over all Feynman diagrams: the compactification absorbs the combinatorial proliferation into the geometry of the boundary strata, and the cohomological nature of the integrand prevents the amplitude from growing factorially.
\end{remark}

\subsubsection{Explicit values through genus 10}

\begin{proposition}[Shadow free energies through genus $10$; \ClaimStatusProvedHere]
% label removed: prop:thqg-I-Fg-values
\index{shadow free energy!explicit values}
Using the Faber--Pandharipande coefficients (Proposition~\ref{prop:thqg-I-faber-pandharipande}):
\begin{align}
F_1(\cA) &= \frac{\kappa}{24}, % label removed: eq:thqg-I-F1\\
F_2(\cA) &= \frac{7\kappa}{5760}, % label removed: eq:thqg-I-F2\\
F_3(\cA) &= \frac{31\kappa}{967680}, % label removed: eq:thqg-I-F3\\
F_4(\cA) &= \frac{127\kappa}{154828800}, % label removed: eq:thqg-I-F4\\
F_5(\cA) &= \frac{73\kappa}{3503554560}, % label removed: eq:thqg-I-F5\\
F_6(\cA) &= \frac{1414477\kappa}{2678117105664000}, % label removed: eq:thqg-I-F6\\
F_7(\cA) &= \frac{8191\kappa}{612141052723200}, % label removed: eq:thqg-I-F7\\
F_8(\cA) &= \frac{16931177\kappa}{49950709902213120000}, % label removed: eq:thqg-I-F8\\
F_9(\cA) &= \frac{5749691557\kappa}{669659197233029971968000}, % label removed: eq:thqg-I-F9\\
F_{10}(\cA) &= \frac{91546277357\kappa}{420928638260761696665600000}. % label removed: eq:thqg-I-F10
\end{align}
\end{proposition}

\begin{proof}
Direct substitution of the Bernoulli numbers into the formula $F_g = \kappa \cdot \frac{2^{2g-1}-1}{2^{2g-1}} \cdot \frac{|B_{2g}|}{(2g)!}$.

At genus $g = 1$: $F_1 = \kappa \cdot \frac{1}{2} \cdot \frac{1/6}{2} = \kappa/24$.

At genus $g = 2$: $(2^3{-}1)/2^3 = 7/8$, $|B_4| = 1/30$, $4! = 24$.  $F_2 = \kappa \cdot \frac{7}{8} \cdot \frac{1/30}{24} = \frac{7\kappa}{5760}$.

At genus $g = 3$: $|B_6| = 1/42$, $6! = 720$, $(2^5 - 1)/2^5 = 31/32$.  $F_3 = \kappa \cdot \frac{31}{32} \cdot \frac{1/42}{720} = \frac{31\kappa}{32 \cdot 30240} = \frac{31\kappa}{967680}$.

At genus $g = 4$: $|B_8| = 1/30$, $8! = 40320$, $(2^7 - 1)/2^7 = 127/128$.  $F_4 = \kappa \cdot \frac{127}{128} \cdot \frac{1/30}{40320} = \frac{127\kappa}{128 \cdot 1209600} = \frac{127\kappa}{154828800}$.

At genus $g = 5$: $|B_{10}| = 5/66$, $10! = 3628800$, $(2^9 - 1)/2^9 = 511/512$.  $F_5 = \kappa \cdot \frac{511}{512} \cdot \frac{5/66}{3628800} = \kappa \cdot \frac{511 \cdot 5}{512 \cdot 66 \cdot 3628800} = \frac{2555\kappa}{122624409600} = \frac{73\kappa}{3503554560}$.

The remaining values are computed similarly.
\end{proof}

\begin{table}[ht]
\centering
\caption{Shadow free energies $F_g(\cA) = \kappa \cdot \lambda_g^{\mathrm{FP}}$ for three families, $g = 1, \ldots, 5$.}
% label removed: tab:thqg-I-Fg-specializations
\index{shadow free energy!table}
\renewcommand{\arraystretch}{1.5}
{\footnotesize
\begin{tabular}{|c|c|c|c|}
\hline
$\boldsymbol{g}$ & $\mathcal{H}_k$, $\kappa = k$ & $\widehat{\mathfrak{sl}}_{2,k}$, $\kappa = \frac{3(k+2)}{4}$ & $\mathrm{Vir}_c$, $\kappa = \frac{c}{2}$ \\
\hline
$1$ & $\dfrac{k}{24}$ & $\dfrac{k+2}{32}$ & $\dfrac{c}{48}$ \\[6pt]
\hline
$2$ & $\dfrac{7k}{5760}$ & $\dfrac{7(k+2)}{7680}$ & $\dfrac{7c}{11520}$ \\[6pt]
\hline
$3$ & $\dfrac{31k}{967680}$ & $\dfrac{31(k+2)}{1290240}$ & $\dfrac{31c}{1935360}$ \\[6pt]
\hline
$4$ & $\dfrac{127k}{154828800}$ & $\dfrac{127(k+2)}{206438400}$ & $\dfrac{127c}{309657600}$ \\[6pt]
\hline
$5$ & $\dfrac{73k}{3503554560}$ & $\dfrac{73(k+2)}{4671406080}$ & $\dfrac{73c}{7007109120}$ \\[6pt]
\hline
\end{tabular}}
\end{table}

\begin{remark}[Numerical magnitudes]
% label removed: rem:thqg-I-numerical-magnitudes
\index{shadow free energy!numerical estimates}
For the Heisenberg at $k = 1$: $F_1 \approx 0.0417$, $F_2 \approx 1.22 \times 10^{-3}$, $F_3 \approx 3.20 \times 10^{-5}$, $F_4 \approx 8.20 \times 10^{-7}$, $F_5 \approx 2.08 \times 10^{-8}$.  The free energies decrease super-exponentially, with ratio $F_{g+1}/F_g \to 1/(2\pi)^2 \approx 0.0253$.  By genus $10$, $F_{10} \approx 2.17 \times 10^{-16}$.  The geometric decay with ratio $1/(4\pi^2) \approx 0.0253$ means that the genus expansion is an extremely rapidly converging series, even at $\hbar = 1$.
\end{remark}

\begin{proposition}[Cumulative partial sums; \ClaimStatusProvedHere]
% label removed: prop:thqg-I-partial-sums
\index{shadow free energy!partial sums}
For the Heisenberg at $k = 1$ and $\hbar = 1$, the partial sums of the genus expansion are:
\begin{align}
S_1 &= F_1 = 0.04167\,, \notag \\
S_2 &= S_1 + F_2 = 0.04288\,, \notag \\
S_3 &= S_2 + F_3 = 0.04291\,, \notag \\
S_4 &= S_3 + F_4 = 0.04291\,, \notag \\
S_5 &= S_4 + F_5 = 0.04291\,, % label removed: eq:thqg-I-partial-sums
\end{align}
converging to $\kappa \cdot (1/\mathrm{sinc}(1/2) - 1) = (1/(2\sin(1/2)) - 1) \cdot 1 \approx 0.04291$.
\end{proposition}

\begin{proof}
Direct numerical evaluation.
\end{proof}


% ======================================================================
%  SUBSECTION 3: CONVERGENCE OF THE GRAVITATIONAL PARTITION FUNCTION
% ======================================================================

\subsection{Convergence of the gravitational partition function}
% label removed: subsec:thqg-I-convergence-grav
\index{gravitational partition function!convergence|textbf}

\subsubsection{Definition of the gravitational partition function}

\begin{definition}[Gravitational partition function]
% label removed: def:thqg-I-grav-partition
\index{gravitational partition function!definition|textbf}
For a holographic modular Koszul datum $\mathcal{H}(T) = (\cA, \cA^!, \mathcal{C}, r(z), \Theta_\cA, \nabla^{\mathrm{hol}})$ and a formal genus-counting parameter~$\hbar$, the \emph{gravitational partition function} is
\begin{equation}% label removed: eq:thqg-I-Z-grav-def
Z_{\mathrm{grav}}(T;\,\hbar) \;:=\; \sum_{g=0}^{\infty} \hbar^g \int_{\overline{\mathcal{M}}_{g}} \operatorname{Sh}_{g,0}(\Theta_\cA)\,.
\end{equation}
The \emph{scalar gravitational partition function} is the restriction to the arity-$0$ scalar channel:
\begin{equation}% label removed: eq:thqg-I-Z-grav-scalar
Z_{\mathrm{grav}}^{\mathrm{scal}}(T;\,\hbar) \;:=\; \sum_{g=0}^{\infty} F_g(\cA)\,\hbar^g \;=\; \kappa(\cA) \cdot \frac{\sqrt{\hbar}/2}{\sin(\sqrt{\hbar}/2)}\,,
\end{equation}
where the closed form follows from Theorem~\ref{thm:thqg-I-generating-function}.
\end{definition}

\begin{remark}[Scalar versus full]
% label removed: rem:thqg-I-scalar-vs-full
The scalar partition function captures only the trace of the modular functor.  The full gravitational partition function includes the higher-arity contributions from the shadow tower: the spectral discriminant $\Delta_\cA$ (arity~$2$ spectral channel), the cubic shadow $\mathfrak{C}$ (arity~$3$), the quartic resonance class $\mathfrak{Q}$ (arity~$4$), and all higher shadows.  For Gaussian algebras (shadow depth $r_{\max} = 2$), the scalar and full partition functions coincide.
\end{remark}

\begin{remark}[Genus-$0$ normalization]
% label removed: rem:thqg-I-genus-0
\index{gravitational partition function!genus-0 normalization}
The genus-$0$ contribution $F_0(\cA) = \kappa(\cA)$ is a convention: it is the ``classical'' contribution from the sphere.  The physical partition function is $Z_{\mathrm{grav}} = \kappa + \sum_{g \geq 1} F_g \hbar^g$, where the leading term $\kappa$ is the tree-level answer and the remaining terms are quantum corrections.  In the semiclassical limit $\hbar \to 0$, only the tree-level term survives.
\end{remark}

\subsubsection{Convergence from HS-sewing and Bernoulli asymptotics}

\begin{theorem}[Perturbative finiteness of twisted gravity; \ClaimStatusProvedHere]
% label removed: thm:thqg-I-perturbative-finiteness
\index{perturbative finiteness!main theorem|textbf}
Let $\cA$ be a modular Koszul chiral algebra satisfying the HS-sewing condition \textup{(}Theorem~\textup{\ref{thm:general-hs-sewing})}.  Then:
\begin{enumerate}[label=\textup{(\roman*)}]
\item \emph{Genus-by-genus finiteness.}  For each $g \geq 1$, the genus-$g$ shadow amplitude $\int_{\overline{\mathcal{M}}_g} \operatorname{Sh}_{g,n}(\Theta_\cA)$ is finite at every arity~$n$.
\item \emph{Scalar partition function.}  The series $\sum_{g=1}^{\infty} F_g(\cA)\,\hbar^g$ converges absolutely for $|\hbar| < 4\pi^2$, with the closed-form expression $Z_{\mathrm{grav}}^{\mathrm{scal}}(\cA;\,\hbar) = \kappa(\cA) \cdot \frac{\sqrt{\hbar}/2}{\sin(\sqrt{\hbar}/2)}$.
\item \emph{Full partition function.}  The full genus-summed partition function $Z_{\mathrm{grav}}(T;\,\hbar)$ converges absolutely for $|\hbar| < R_{\mathrm{grav}}$, where $R_{\mathrm{grav}} \geq 4\pi^2$.  For Gaussian algebras, $R_{\mathrm{grav}} = \infty$.
\item \emph{No UV renormalization.}  No counterterms or renormalization are needed at any genus.  The Maurer--Cartan equation $D_\cA\Theta_\cA + \tfrac{1}{2}[\Theta_\cA,\Theta_\cA] = 0$ (Theorem~\ref*{V1-thm:mc2-bar-intrinsic}) is an exact identity, and the Fulton--MacPherson compactification provides the geometric regulator.
\end{enumerate}
\end{theorem}

\begin{proof}
\emph{Part~(i).}  Fix genus~$g$ and arity~$n$.  The shadow representation $\operatorname{Sh}_{g,n}(\Theta_\cA) \in H^*(\overline{\mathcal{M}}_{g,n})$ is a cohomology class on a compact orbifold.  Integration over $\overline{\mathcal{M}}_{g,n}$ is a continuous linear functional on this finite-dimensional space, hence bounded.  The nontrivial content is that the class itself is well defined: this follows from the bar-intrinsic construction of $\Theta_\cA$ (Theorem~\ref*{V1-thm:mc2-bar-intrinsic}), which gives $\Theta_\cA := D_\cA - d_0 \in \mathrm{MC}(\gAmod)$ without convergence issues.

At the chain level, the integral is computed by the arity-cutoff lemma (Lemma~\ref{lem:arity-cutoff}): the graph sum through arity~$r$ at genus~$g$ involves only finitely many graphs $\Gamma$ with $|\Gamma| \leq r$ edges and loop genus $\leq g$.  Each graph amplitude $\ell_\Gamma = |\Aut(\Gamma)|^{-1} \int_{\overline{\mathcal{M}}_{g(\Gamma),n(\Gamma)}} \omega_\Gamma$ is a period integral on a compact moduli space, hence finite.

\emph{Part~(ii).}  This is Theorem~\ref{thm:thqg-I-absolute-convergence}: the Bernoulli asymptotics $|F_g| \sim 2|\kappa|/(2\pi)^{2g}$ give absolute convergence for $|\hbar| < 4\pi^2$, and analytic continuation via the closed form extends to all of~$\mathbb{C}$ as a meromorphic function.

\emph{Part~(iii).}  For the full partition function, each genus-$g$ contribution is a sum over the shadow tower.  The shadow tower truncation $\Theta_\cA^{\leq r}$ is a finite sum at each arity, and the inverse limit $\Theta_\cA = \varprojlim \Theta_\cA^{\leq r}$ converges (Theorem~\ref{thm:recursive-existence}).  The arity-$r$ contribution to $F_g$ is bounded by $C_r / (2\pi)^{2g}$, so the genus sum converges at each arity.  For Gaussian algebras ($r_{\max} = 2$), all arity-$r \geq 3$ contributions vanish, and the full partition function equals the scalar partition function; the latter converges for all $|\hbar| < 4\pi^2$, and the higher-arity vanishing means $R_{\mathrm{grav}} = \infty$ (the function is entire).

\emph{Part~(iv).}  The absence of UV divergences is the combined consequence of three features:
\begin{enumerate}[label=\textup{(\alph*)}]
\item \emph{Compactness.}  The Fulton--MacPherson compactification $\overline{\mathrm{FM}}_n(\Sigma_g)$ is a compact manifold with corners.  All integrals are over compact domains.
\item \emph{Exactness.}  The Maurer--Cartan equation is an identity of the bar differential ($D_\cA^2 = 0$), not a perturbative expansion.  There are no UV divergences to regulate because the equation is exact, not asymptotic.
\item \emph{Arity cutoff.}  At each genus~$g$, only finitely many graph topologies contribute (Lemma~\ref{lem:arity-cutoff}).  The number of contributing graphs at genus $g$ and arity $r$ is bounded by a combinatorial function of $(g,r)$, and each graph amplitude is a finite integral.
\end{enumerate}
\end{proof}

\subsubsection{Comparison with standard QFT divergences}

\begin{remark}[Contrast with perturbative quantum gravity in $d = 4$]
% label removed: rem:thqg-I-contrast-4d
\index{UV divergence!standard quantum gravity}
In standard four-dimensional perturbative quantum gravity (Einstein--Hilbert action), the genus-$g$ contribution to the graviton scattering amplitude diverges for $g \geq 2$ as
\begin{equation}% label removed: eq:thqg-I-4d-divergence
\mathcal{A}_g \;\sim\; \Lambda_{\mathrm{UV}}^{2(g-1)} \int_{\mathcal{M}_g} |\Omega_g|^2 \cdot R^{2g},
\end{equation}
where $\Lambda_{\mathrm{UV}}$ is the UV cutoff and $R$ is the Riemann curvature.  In the twisted holographic setting, the analogous quantity $\mathcal{A}_g^{\mathrm{tw}} = \int_{\overline{\mathcal{M}}_g} \operatorname{Sh}_{g,n}(\Theta_\cA)$ is finite at every genus because: (a)~the integration domain $\overline{\mathcal{M}}_g$ is compact; (b)~the integrand is a smooth cohomology class; (c)~the Maurer--Cartan equation provides an exact Ward identity constraining the amplitudes genus by genus.  The fundamental difference is that the twisted holographic gravity lives on the moduli space of curves, not on spacetime.
\end{remark}

\begin{remark}[Contrast with string theory]
% label removed: rem:thqg-I-contrast-string
\index{string theory!UV finiteness}
In string theory, genus-$g$ amplitudes are finite individually (after integration over $\overline{\mathcal{M}}_g$), but the genus expansion $\sum_g g_s^{2g-2} \mathcal{A}_g$ is asymptotic: it diverges factorially as $\mathcal{A}_g \sim (2g)! \cdot \alpha'^{-g}$.  In the twisted holographic setting, the individual amplitudes are also finite (same reason: compactness of $\overline{\mathcal{M}}_g$), but the genus expansion \emph{converges}: $F_g \sim \kappa/(2\pi)^{2g}$ decays exponentially.  The difference stems from the cohomological nature of the shadow amplitudes: in string theory, the integrand over $\mathcal{M}_g$ is a genuine function that grows with $g$; in the shadow theory, the integrand is a fixed cohomology class $\kappa \cdot \lambda_g$ whose integral $\lambda_g^{\mathrm{FP}}$ decays.
\end{remark}

\subsubsection{The role of Koszulness in finiteness}

\begin{proposition}[Koszulness and finiteness; \ClaimStatusProvedHere]
% label removed: prop:thqg-I-koszulness-finiteness
\index{Koszulness!role in finiteness}
For a modular Koszul chiral algebra $\cA$, the shadow tower $\Theta_\cA^{\leq r}$ has the following finiteness properties:
\begin{enumerate}[label=\textup{(\roman*)}]
\item At each arity $r$ and genus $g$, the graph sum
  \begin{equation}% label removed: eq:thqg-I-graph-sum-finite
  \Theta_\cA^{(g,r)} \;=\; \sum_{\substack{\Gamma \in \Gmod(g,r) \\ \Gamma \text{ connected}}} \frac{1}{|\Aut(\Gamma)|} \int_{\overline{\mathcal{M}}_{g(\Gamma),n(\Gamma)}} \omega_\Gamma
  \end{equation}
  is a finite sum over a finite set of graphs.
\item The Koszul property ensures that the transferred cyclic minimal model has only finitely many nontrivial operations at each arity, so the vertex labels $\ell_n^{\mathrm{tr}}$ are nonzero only for $n$ in a finite range.
\item The shadow depth $r_{\max}(\cA)$ bounds the arity of contributing graphs: for $r > r_{\max}$, the obstruction class $o_{r+1}(\cA) = 0$ and the shadow tower stabilizes.
\end{enumerate}
\end{proposition}

\begin{proof}
Part~(i) is the arity-cutoff lemma (Lemma~\ref{lem:arity-cutoff}).  The set $\Gmod(g,r)$ of connected modular graphs with loop genus $g$ and total arity $\leq r$ is finite: the number of vertices $v$, edges $e$, and external legs $n$ satisfy $v - e + n = 1 - g$ (the Euler characteristic), $2e + n \leq rv$ (each vertex has at most $r$ half-edges), and $g = e - v + 1 + \sum_{v'} g_{v'}$ where $g_{v'}$ is the genus decoration at each vertex.  These constraints bound $v, e, n$ in terms of $g$ and $r$.

Part~(ii) follows from the PBW concentration property (Proposition~\ref{prop:standard-examples-modular-koszul}): for a Koszul algebra, the transferred $A_\infty$ operations $m_n^{\mathrm{tr}}$ vanish for $n$ sufficiently large (bounded by the Koszul dual degree).

Part~(iii) is the definition of shadow depth: if $o_{r+1}(\cA) = 0$, then $\Theta_\cA^{\leq r+1} = \Theta_\cA^{\leq r}$, and the tower stabilizes.
\end{proof}

\begin{proposition}[Effective graph count; \ClaimStatusProvedHere]
% label removed: prop:thqg-I-graph-count
\index{graph sum!effective count}
The number of connected modular graphs contributing to $\Theta_\cA^{(g,r)}$ is bounded by
\begin{equation}% label removed: eq:thqg-I-graph-bound
|\Gmod(g,r)| \;\leq\; C(g,r) \;:=\; \sum_{v=1}^{3g-3+r} \binom{v(r-1)/2 + g - 1}{g} \cdot r^v,
\end{equation}
which is finite for fixed $(g,r)$ and grows polynomially in $r$ at fixed $g$.
\end{proposition}

\begin{proof}
A connected graph with $v$ vertices and loop genus $g$ has $e = v + g - 1$ edges.  Each vertex has valence at most $r$ (by the arity cutoff), so $2e + n \leq rv$, giving $n \leq rv - 2(v+g-1) = v(r-2) - 2g + 2$.  The condition $n \geq 0$ requires $v \geq (2g-2)/(r-2)$ (for $r \geq 3$).  The number of graphs with $v$ labeled vertices and $e$ edges is at most $\binom{v(v-1)/2}{e}$ (choosing edges from all possible pairs); the symmetry group reduces this, but the crude bound suffices.  The factor $r^v$ accounts for the vertex labeling (choice of operation at each vertex from $\{m_2, m_3, \ldots, m_r\}$).
\end{proof}

\subsubsection{The structure of finiteness at each genus}

\begin{proposition}[Genus-$g$ amplitude decomposition; \ClaimStatusProvedHere]
% label removed: prop:thqg-I-genus-g-decomposition
\index{genus-$g$ amplitude!decomposition}
The genus-$g$ shadow amplitude decomposes as
\begin{equation}% label removed: eq:thqg-I-genus-g-decomp
\int_{\overline{\mathcal{M}}_g} \operatorname{Sh}_{g,0}(\Theta_\cA) \;=\; F_g^{\mathrm{scal}}(\cA) \;+\; \sum_{r=2}^{r_{\max}} F_g^{(r)}(\cA),
\end{equation}
where $F_g^{\mathrm{scal}} = \kappa \cdot \lambda_g^{\mathrm{FP}}$ is the scalar (arity-$0$) contribution and $F_g^{(r)}$ is the arity-$r$ correction.  For a Gaussian algebra ($r_{\max} = 2$), the sum is empty and only the scalar term survives.  For a Lie-type algebra ($r_{\max} = 3$), there is a single cubic correction $F_g^{(3)}$.  For a contact-type algebra ($r_{\max} = 4$), there are cubic and quartic corrections.  For a mixed algebra ($r_{\max} = \infty$), the sum is infinite but converges by the shadow tower convergence theorem (Theorem~\ref{thm:recursive-existence}).
\end{proposition}

\begin{proof}
The shadow tower decomposes by arity: $\Theta_\cA = \sum_{r \geq 0} \Theta_\cA^{(r)}$ where $\Theta_\cA^{(r)}$ is the arity-$r$ component.  The genus-$g$ projection at arity $0$ gives $F_g^{\mathrm{scal}}$ by the universality theorem.  At arity $r \geq 2$, the projection $\int_{\overline{\mathcal{M}}_g} \operatorname{Sh}_{g,r}(\Theta_\cA)$ involves a graph sum with vertex labels from the arity-$r$ operations, and each term is finite by the arity-cutoff argument.  The vanishing of $F_g^{(r)}$ for $r > r_{\max}$ follows from the shadow depth.
\end{proof}


% ======================================================================
%  SUBSECTION 4: THE HEISENBERG PARTITION FUNCTION AS PROTOTYPE
% ======================================================================

\subsection{The Heisenberg partition function as prototype}
% label removed: subsec:thqg-I-heisenberg-prototype
\index{Heisenberg!partition function|textbf}

The Heisenberg algebra is the archetype: Gaussian shadow depth ($r_{\max} = 2$), exact solvability via second quantization, and explicit Fredholm determinant formulas at all genera.

\subsubsection{The Heisenberg at genus $g$}

\begin{theorem}[Heisenberg gravitational partition function; \ClaimStatusProvedElsewhere{} {\textup{Theorem~\ref{thm:heisenberg-sewing}}}]
% label removed: thm:thqg-I-heisenberg-partition
\index{Heisenberg!genus-$g$ partition function}
For the rank-$1$ Heisenberg algebra $\mathcal{H}_k$ on a genus-$g$ surface $\Sigma_g$ with period matrix $\Omega \in \mathbb{H}_g$ \textup{(}the Siegel upper half-space\textup{)}, the genus-$g$ partition function is the Fredholm determinant
\begin{equation}% label removed: eq:thqg-I-heisenberg-Zg
\boxed{ Z_g(\mathcal{H}_k;\,\Omega) \;=\; \bigl(\det\operatorname{Im}\Omega\bigr)^{-k/2} \cdot \bigl(\det{}'_\zeta \Delta_{\Sigma_g}\bigr)^{-k/2}, }
\end{equation}
where $\det{}'_\zeta$ denotes the zeta-regularized determinant of the scalar Laplacian $\Delta_{\Sigma_g}$ on $\Sigma_g$ \textup{(}excluding the zero mode\textup{)}.  The shadow free energy extracted from this is
\begin{equation}% label removed: eq:thqg-I-heisenberg-Fg
F_g(\mathcal{H}_k) \;=\; k \cdot \frac{2^{2g-1}-1}{2^{2g-1}} \cdot \frac{|B_{2g}|}{(2g)!}\,,
\end{equation}
confirming the general formula~\eqref{eq:thqg-I-Fg-formula} at $\kappa(\mathcal{H}_k) = k$.
\end{theorem}

\begin{proof}
By Wick's theorem, all amplitudes of the Heisenberg algebra factorize into propagator pairings.  The propagator on a genus-$g$ surface is the Green's function
\begin{equation}% label removed: eq:thqg-I-green-genus-g
G(z,w;\,\Omega) \;=\; \log|E(z,w;\,\Omega)| - \frac{(\operatorname{Im}\int_w^z \omega)^T (\operatorname{Im}\Omega)^{-1} (\operatorname{Im}\int_w^z \omega)}{2},
\end{equation}
where $E(z,w;\,\Omega)$ is the prime form and $\omega = (\omega_1, \ldots, \omega_g)^T$ is the normalized basis of holomorphic differentials.  The prime form is a section of $K^{-1/2} \boxtimes K^{-1/2}$ (not $K^{+1/2}$).

The partition function is the functional determinant of the kinetic operator:
\[
Z_g \;=\; (\det{}'_\zeta \bar{\partial}^*\bar{\partial})^{-k/2} \;=\; (\det\operatorname{Im}\Omega)^{-k/2}\, (\det{}'_\zeta \Delta)^{-k/2},
\]
where the factorization into period-matrix and Laplacian contributions is the Belavin--Knizhnik factorization~\cite{BK86}.  The Polyakov formula identifies $\det{}'_\zeta \Delta$ with the norm of the Mumford form on $\overline{\mathcal{M}}_g$; its $g$-th Chern class is $\lambda_g$.  The shadow free energy is obtained by integrating the Chern class:
\[
F_g(\mathcal{H}_k) = k \int_{\overline{\mathcal{M}}_g} \lambda_g = k \cdot \lambda_g^{\mathrm{FP}}.
\]
\end{proof}

\begin{remark}[The Belavin--Knizhnik factorization]
% label removed: rem:thqg-I-belavin-knizhnik
\index{Belavin--Knizhnik!factorization}
The factorization $\det{}'_\zeta(\bar\partial^*\bar\partial) = \det(\operatorname{Im}\Omega) \cdot \det{}'_\zeta(\Delta)$ separates the ``holomorphic'' contribution (the period matrix, which depends on the complex structure of $\Sigma_g$) from the ``analytic'' contribution (the Laplacian determinant, which depends on the conformal class).  The period matrix factor $(\det\operatorname{Im}\Omega)^{-k/2}$ is the modular-covariant piece that transforms under $\mathrm{Sp}(2g,\mathbb{Z})$, while the Laplacian factor is modular-invariant.  Together they give the conformal block for the Heisenberg algebra on $\Sigma_g$.
\end{remark}

\subsubsection{The Fredholm determinant formula}

\begin{proposition}[Fredholm determinant representation; \ClaimStatusProvedElsewhere{} {\textup{Theorem~\ref{thm:heisenberg-one-particle-sewing}}}]
% label removed: prop:thqg-I-fredholm
\index{Fredholm determinant!Heisenberg|textbf}
The genus-$g$ partition function admits a Fredholm determinant representation:
\begin{equation}% label removed: eq:thqg-I-fredholm-det
Z_g(\mathcal{H}_k;\,\Omega) \;=\; \det(1 - K_g)^{-k},
\end{equation}
where $K_g$ is the trace-class sewing operator on the one-particle Bergman space $A^2(\partial D)$, with eigenvalues $q_1^{n_1} \cdots q_{3g-3}^{n_{3g-3}}$ indexed by the sewing moduli $(q_1, \ldots, q_{3g-3})$.  The trace-class norm satisfies
\begin{equation}% label removed: eq:thqg-I-Kg-trace-norm
\|K_g\|_1 \;=\; \prod_{i=1}^{3g-3} \frac{q_i}{1-q_i} \;<\; \infty \quad\text{for all } q_i \in (0,1).
\end{equation}
\end{proposition}

\begin{proof}
The pair-of-pants amplitude for the Heisenberg is the second quantization $\Gamma(R)$ of the one-particle Bergman restriction $R \colon A^2(D_{\mathrm{out}}) \to A^2(D_{\mathrm{in},1}) \otimes A^2(D_{\mathrm{in},2})$ (Theorem~\ref{thm:heisenberg-one-particle-sewing}).  The one-particle restriction has matrix entries $R_{n,m} = \frac{k}{2\pi i} \oint\oint z^{-n-1}w^{-m-1}G_{\mathrm{pants}}(z,w)\,dz\,dw$, satisfying $|R_{n,m}| \leq C\,q^{(n+m)/2}$ from the collar length $t = -\log q$.  Hence $R$ is Hilbert--Schmidt and $\Gamma(R)$ is trace class.

The composition of $3g - 3$ sewing operators gives the genus-$g$ operator $K_g = R_1 \circ \cdots \circ R_{3g-3}$ on the one-particle space.  The Fredholm determinant identity $\Tr\,\Gamma(K_g)^k = \det(1 - K_g)^{-k}$ is the boson-fermion correspondence.  The trace-class norm bound follows from $\|K_g\|_1 \leq \prod_{i=1}^{3g-3} \|R_i\|_1 = \prod_{i=1}^{3g-3} q_i/(1-q_i)$.
\end{proof}

\begin{remark}[The Bergman space]
% label removed: rem:thqg-I-bergman
\index{Bergman space}
The one-particle Bergman space $A^2(\partial D)$ consists of holomorphic functions on the annular collar around a sewing circle, square-integrable with respect to the Bergman kernel.  An orthonormal basis is $\{z^n / \sqrt{2\pi}\}_{n \geq 1}$, and the one-particle sewing operator in this basis has matrix elements $K_{nm} = \delta_{nm} q^n$.  The Fredholm determinant is therefore $\det(1 - K) = \prod_{n=1}^\infty (1 - q^n)$, giving the Dedekind eta function upon including the zero-point energy.
\end{remark}

\subsubsection{Comparison with the Dedekind eta function}

\begin{example}[Genus $1$: the Dedekind eta function]
% label removed: ex:thqg-I-genus1-eta
\index{Dedekind eta!Heisenberg partition function}
At genus $g = 1$, the period matrix is $\Omega = \tau \in \mathbb{H}$ and the sewing parameter is $q = e^{2\pi i\tau}$.  The Fredholm determinant becomes
\begin{equation}% label removed: eq:thqg-I-genus1-fredholm
Z_1(\mathcal{H}_k;\,\tau) \;=\; (\operatorname{Im}\tau)^{-k/2}\, \prod_{n=1}^{\infty} (1 - q^n)^{-k} \;=\; (\operatorname{Im}\tau)^{-k/2}\, \bigl|\eta(\tau)\bigr|^{-2k},
\end{equation}
where $\eta(\tau) = q^{1/24}\prod_{n=1}^{\infty}(1 - q^n)$ is the Dedekind eta function.  The shadow free energy at genus $1$ is $F_1(\mathcal{H}_k) = k/24$, matching~\eqref{eq:thqg-I-F1} with $\kappa = k$.

The connection to the modular characteristic is transparent: the $q$-expansion of $\log \eta(\tau) = \frac{2\pi i\tau}{24} - \sum_{n=1}^\infty \log(1-q^n)$ begins with $\frac{2\pi i \tau}{24}$, and the genus-$1$ free energy is the constant term of the $q$-expansion of $-k\log|\eta(\tau)|$, which is $k/24$.
\end{example}

\begin{example}[Genus $2$: explicit formula]
% label removed: ex:thqg-I-genus2
\index{Heisenberg!genus-2 partition function}
At genus $g = 2$, the period matrix $\Omega = \begin{psmallmatrix} \tau_1 & \tau_{12} \\ \tau_{12} & \tau_2 \end{psmallmatrix} \in \mathbb{H}_2$ and the partition function is
\begin{equation}% label removed: eq:thqg-I-genus2
Z_2(\mathcal{H}_k;\,\Omega) \;=\; \bigl(\det\operatorname{Im}\Omega\bigr)^{-k/2} \cdot \bigl(\det{}'_\zeta \Delta_{\Sigma_2}\bigr)^{-k/2}.
\end{equation}
The shadow free energy is $F_2(\mathcal{H}_k) = 7k/5760$, confirming~\eqref{eq:thqg-I-F2}.

The Fredholm determinant at genus $2$ involves the sewing of three pairs of pants along three circles.  In the Schottky uniformization, $\Sigma_2 = (\mathbb{P}^1 \setminus \{6 \text{ discs}\}) / \langle \gamma_1, \gamma_2, \gamma_3 \rangle$ with three sewing parameters $q_1, q_2, q_3$.  The one-particle sewing operator is $K_2 = R_1 \circ R_2 \circ R_3$, and $\det(1 - K_2) = \prod_{n_1,n_2,n_3 \geq 1} (1 - q_1^{n_1} q_2^{n_2} q_3^{n_3})$.  The shadow free energy is obtained by integrating over the Siegel modular variety $\mathcal{M}_2 = \mathbb{H}_2 / \mathrm{Sp}(4,\mathbb{Z})$.
\end{example}

\begin{example}[Genus $3$]
% label removed: ex:thqg-I-genus3-heisenberg
\index{Heisenberg!genus-3 partition function}
At genus $g = 3$, the surface $\Sigma_3$ has $3g - 3 = 6$ sewing parameters.  The Fredholm determinant is $\det(1 - K_3)^{-k}$ where $K_3$ is a trace-class operator on the one-particle Bergman space.  The shadow free energy is $F_3(\mathcal{H}_k) = 31k/967680$.  Numerically at $k = 1$: $F_3 \approx 3.20 \times 10^{-5}$, which is already three orders of magnitude smaller than $F_1 = 1/24 \approx 0.042$.
\end{example}

\subsubsection{One-particle Bergman reduction}

\begin{remark}[One-particle reduction principle]
% label removed: rem:thqg-I-one-particle
\index{Bergman space!one-particle reduction|textbf}
The Heisenberg partition function factorizes completely into one-particle data; this is the second quantization principle.  The one-particle sewing operator $T_q$ on the Bergman space $A^2(\partial D)$ has eigenvalues $q^n$ ($n = 1, 2, 3, \ldots$), and the full partition function is
\[
Z_g(\mathcal{H}_k) \;=\; \det(1 - T_{q_1} \circ \cdots \circ T_{q_{3g-3}})^{-k} \;=\; \exp\!\left( k \sum_{m=1}^{\infty} \frac{1}{m} \Tr\bigl(T_{q_1} \cdots T_{q_{3g-3}}\bigr)^m \right).
\]
This exponential formula makes the connection to the shadow free energy transparent: $F_g$ is the constant term of $\log Z_g$ in the expansion in sewing moduli.  The one-particle reduction is specific to the Heisenberg (Gaussian shadow depth $r_{\max} = 2$): for non-Gaussian algebras, the multi-particle sector contributes genuinely new data that cannot be recovered from the one-particle space.
\end{remark}

\begin{proposition}[Exponential representation; \ClaimStatusProvedHere]
% label removed: prop:thqg-I-exponential-rep
\index{Fredholm determinant!exponential representation}
The Fredholm determinant admits the exponential representation
\begin{equation}% label removed: eq:thqg-I-exponential-fredholm
\log \det(1 - K_g)^{-k} \;=\; k \sum_{m=1}^\infty \frac{1}{m} \Tr(K_g^m) \;=\; k \sum_{m=1}^\infty \frac{1}{m} \sum_{\mathbf{n} \in \mathbb{Z}_{\geq 1}^{3g-3}} \prod_{i=1}^{3g-3} q_i^{m n_i},
\end{equation}
which converges absolutely for all $q_i \in (0,1)$.  The shadow free energy $F_g$ is extracted as the coefficient of the leading term in the $\mathbf{q}$-expansion, integrated over $\overline{\mathcal{M}}_g$ with the Weil--Petersson measure.
\end{proposition}

\begin{proof}
The Fredholm determinant identity $\log\det(1 - K) = \Tr\log(1 - K) = -\sum_{m=1}^\infty \frac{1}{m}\Tr(K^m)$ holds for trace-class $K$ with $\|K\|_1 < 1$.  For $q_i \in (0,1)$, the eigenvalues of $K_g$ are $\prod_{i=1}^{3g-3} q_i^{n_i}$ with $n_i \geq 1$, all less than $1$.  The trace $\Tr(K_g^m) = \sum_{\mathbf{n}} \prod q_i^{mn_i}$ converges absolutely because $\sum_{\mathbf{n}} \prod q_i^{n_i} = \prod_{i=1}^{3g-3} q_i/(1-q_i) < \infty$.
\end{proof}

\begin{remark}[Rank-$d$ Heisenberg]
% label removed: rem:thqg-I-rank-d-heisenberg
\index{Heisenberg!rank-$d$}
For the rank-$d$ Heisenberg algebra $\mathcal{H}_k^{\otimes d}$ (i.e., $d$ independent free bosons at level $k$), the partition function is the $d$-th power of the rank-$1$ answer:
\begin{equation}% label removed: eq:thqg-I-rank-d
Z_g(\mathcal{H}_k^{\otimes d};\,\Omega) \;=\; Z_g(\mathcal{H}_k;\,\Omega)^d \;=\; \bigl(\det\operatorname{Im}\Omega\bigr)^{-dk/2} \cdot \bigl(\det{}'_\zeta\Delta_{\Sigma_g}\bigr)^{-dk/2}.
\end{equation}
The modular characteristic is $\kappa(\mathcal{H}_k^{\otimes d}) = dk$, and the shadow free energy is $F_g(\mathcal{H}_k^{\otimes d}) = dk \cdot \lambda_g^{\mathrm{FP}}$.  This is consistent with the independent sum factorization (Proposition~\ref{prop:independent-sum-factorization}): for $\cA = \cA_1 \oplus \cA_2$ with vanishing mixed OPE, $\kappa(\cA) = \kappa(\cA_1) + \kappa(\cA_2)$.
\end{remark}


% ======================================================================
%  SUBSECTION 5: FINITENESS FOR THE STANDARD LANDSCAPE
% ======================================================================

\subsection{Finiteness for the standard landscape}
% label removed: subsec:thqg-I-standard-landscape
\index{standard landscape!finiteness|textbf}

\subsubsection{Verification of hypotheses (H1)--(H2)}

The HS-sewing criterion (Theorem~\ref{thm:general-hs-sewing}) requires two conditions:
\begin{enumerate}[label=\textbf{(H\arabic*)}]
\item % label removed: item:thqg-I-H1 \emph{Subexponential sector growth:} $\log \dim H_n = o(n)$.
\item % label removed: item:thqg-I-H2 \emph{Polynomial OPE growth:} $|C^{c,k}_{a,i;\,b,j}| \leq K(a+b+c+1)^N$.
\end{enumerate}

\begin{proposition}[Standard families satisfy HS-sewing; \ClaimStatusProvedElsewhere{} {\textup{Theorem~\ref{thm:general-hs-sewing}}}]
% label removed: prop:thqg-I-standard-hs
\index{HS-sewing!standard families}
Every algebra in the following list satisfies \ref{item:thqg-I-H1} and \ref{item:thqg-I-H2}:
\begin{enumerate}[label=\textup{(\alph*)}]
\item Heisenberg $\mathcal{H}_k$ (any rank, any level $k \neq 0$).
\item Affine Kac--Moody $\widehat{\fg}_k$ ($\fg$ simple, $k \neq -h^\vee$).
\item Virasoro $\mathrm{Vir}_c$ (any $c \in \mathbb{C}$).
\item $\beta\gamma$ system (any central charge).
\item Principal $\mathcal{W}$-algebras $\mathcal{W}_k(\fg)$ ($\fg$ simple, generic $k$).
\item Lattice VOAs $V_\Lambda$ ($\Lambda$ even integral).
\item Free fermions $\psi_1, \ldots, \psi_n$ (any rank).
\end{enumerate}
\end{proposition}

\begin{proof}
For each family, we verify \ref{item:thqg-I-H1} and \ref{item:thqg-I-H2} by direct computation.  The key estimates were computed in Computations~\ref{comp:thqg-I-hs-heisenberg}--\ref{comp:thqg-I-hs-fermion}.

\emph{(a) Heisenberg.}  $\dim H_n = p(n)$, $\log p(n) \sim \pi\sqrt{2n/3} = o(n)$.  OPE degree $N = 1$.

\emph{(b) Affine Kac--Moody.}  $\dim H_n$ has subexponential growth (Kac--Peterson~\cite{KP84}).  OPE degree $N = 1$.

\emph{(c) Virasoro.}  $\dim H_n = p(n)$, $\log p(n) = o(n)$.  OPE degree $N = 2$.

\emph{(d) $\beta\gamma$.}  $\dim H_n = p(n)^2$, $\log \dim H_n = o(n)$.  OPE degree $N = 1$.

\emph{(e) $\mathcal{W}_N$.}  Sector growth subexponential with exponent $O(n^{(N-1)/N})$.  OPE degree $N_{\max} = 2N$.

\emph{(f) Lattice VOA.}  $\dim H_n = |\Lambda^*/\Lambda|\cdot p_r(n)$; $\log p_r(n) = O(\sqrt{n})$.  OPE degree $\leq r$.

\emph{(g) Free fermions.}  $\dim H_n$ grows at most as $2^n$ (number of subsets of $n$ fermion modes), but the conformal-weight grading gives polynomial growth.  OPE degree $N = 0$.
\end{proof}

\subsubsection{Detailed verification for each family}

\begin{proposition}[Heisenberg: detailed HS verification; \ClaimStatusProvedHere]
% label removed: prop:thqg-I-heisenberg-detail
\index{Heisenberg!HS verification}
For the rank-$1$ Heisenberg $\mathcal{H}_k$ at level $k \neq 0$:
\begin{enumerate}[label=\textup{(\roman*)}]
\item The OPE structure constants are $C^{n-1,\lambda}_{n,\mu;\,0,\mathrm{vac}} = k \cdot \delta_{|\lambda|-|\mu|,1}$ for the fundamental mode $a_{-1}|0\rangle$ acting on a partition state $|\mu\rangle$ of weight $n$ to produce a partition state $|\lambda\rangle$ of weight $n-1$.  The general OPE coefficient at weight level $(a,b,c)$ with $a+b-c = m$ (the number of contractions) satisfies $|C^{c,\lambda}_{a,\mu;\,b,\nu}| \leq |k|^m \cdot m! \leq |k|^{a+b} \cdot (a+b)!$, which is polynomial for fixed $m$.
\item The HS norm at weight level $n$ satisfies the exact formula
  \begin{equation}% label removed: eq:thqg-I-heisenberg-exact-hs
  \sum_{\substack{a+b+c = n}} \|m^c_{a,b}\|_{\HS}^2 \;=\; |k|^2 \cdot \sum_{\substack{a+b+c=n}} p(a)p(b)p(c) \cdot \#\{\text{Wick contractions}\},
  \end{equation}
  where the number of Wick contractions at level $n$ grows polynomially.
\end{enumerate}
\end{proposition}

\begin{proof}
The Heisenberg OPE is computed by Wick's theorem: each pair of mode operators $a_m, a_n$ contracts to $k \cdot m \cdot \delta_{m+n,0}$.  The number of contractions between a state $|\mu\rangle$ of weight $a$ and a state $|\nu\rangle$ of weight $b$ is at most $\min(a,b)$, and each contraction produces a factor of $|k| \cdot n$ for a mode of level $n$.  The polynomial bound follows from the multinomial structure of the Wick expansion.
\end{proof}

\begin{proposition}[Affine Kac--Moody: detailed HS verification; \ClaimStatusProvedHere]
% label removed: prop:thqg-I-affine-detail
\index{affine Kac--Moody!HS verification}
For $\widehat{\fg}_k$ with $\fg$ simple of dimension $d$, rank $r$, and dual Coxeter number $h^\vee$, at non-critical level $k \neq -h^\vee$:
\begin{enumerate}[label=\textup{(\roman*)}]
\item The Kac--Peterson character formula gives $\chi_k(\tau) = \sum_{n \geq 0} (\dim H_n) q^n = q^{-c_{\text{eff}}/24} \cdot \Theta_{k+h^\vee,\fg}(\tau) / \eta(\tau)^d$, where $\Theta_{k+h^\vee,\fg}$ is the affine theta function.  The dimension $\dim H_n$ is dominated by $\eta(\tau)^{-d}$, giving $\log \dim H_n \sim \pi\sqrt{2dn/3}$.
\item The OPE of currents $J^a_m J^b_n = k m \delta^{ab}\delta_{m+n,0} + f^{ab}{}_c J^c_{m+n}$ has structure constants bounded by $\max(|k|, \|f\|_\infty) \cdot n$ for mode number $n$, which is polynomial of degree $N = 1$.
\item The HS sewing sum satisfies
  \begin{equation}% label removed: eq:thqg-I-affine-hs-refined
  \sum_{a,b,c \geq 0} q^{a+b+c} \|m^c_{a,b}\|_{\HS}^2 \;\leq\; (|k| + \|f\|_\infty)^2 \cdot \eta(q)^{-2d} \cdot P(q),
  \end{equation}
  where $P(q) = \sum_{n \geq 0} n^2 (\dim H_n)^2 q^n$ is a convergent power series for $q \in (0,1)$.
\end{enumerate}
\end{proposition}

\begin{proof}
(i) is the Kac--Peterson asymptotic~\cite{KP84}.  (ii) follows from the affine Lie algebra commutation relations.  (iii) combines the estimates from (i) and (ii) with the general HS bound~\eqref{eq:thqg-I-hs-criterion-bound}.
\end{proof}

\begin{proposition}[Virasoro: detailed HS verification; \ClaimStatusProvedHere]
% label removed: prop:thqg-I-virasoro-detail
\index{Virasoro!HS verification}
For $\mathrm{Vir}_c$ at central charge $c$:
\begin{enumerate}[label=\textup{(\roman*)}]
\item The Verma module at weight $h$ has $\dim H_n = p(n)$ (partitions of $n$, corresponding to states $L_{-n_1}L_{-n_2}\cdots L_{-n_k}|h\rangle$ with $n_1 + \cdots + n_k = n$).
\item The OPE $T(z)T(w) \sim c/(2(z-w)^4) + 2T(w)/(z-w)^2 + \partial T(w)/(z-w)$ has a quartic pole.  The resulting structure constants satisfy $|C^c_{a,b}| \leq (|c|/2 + 2) \cdot (a+b+1)^2$, which is polynomial of degree $N = 2$.
\item The refined HS bound:
  \begin{equation}% label removed: eq:thqg-I-virasoro-refined-hs
  \sum_{a,b,c \geq 0} q^{a+b+c} \|m^c_{a,b}\|_{\HS}^2 \;\leq\; (|c|/2 + 2)^2 \sum_{n \geq 0} n^4 p(n)^3 q^n.
  \end{equation}
  The series $\sum n^4 p(n)^3 q^n$ converges for all $q \in (0,1)$ because $p(n)^3 \leq e^{3\pi\sqrt{2n/3} + O(\log n)}$ grows subexponentially while $q^n$ decays exponentially.  Concretely, for $q = 0.5$: the $n = 100$ term is $100^4 \cdot p(100)^3 \cdot 0.5^{100} \approx 10^8 \cdot (1.9 \times 10^8)^3 \cdot 10^{-30} \approx 10^{2}$, and the series converges.
\end{enumerate}
\end{proposition}

\begin{proof}
(i) is standard Virasoro representation theory.  (ii) follows from the mode algebra $[L_m, L_n] = (m-n)L_{m+n} + \frac{c}{12}m(m^2-1)\delta_{m+n,0}$, which gives $|C| \leq |c|/12 \cdot m^3 + |m-n| \cdot |L_{m+n}|$; the polynomial bound is obtained by tracking the maximal power of conformal weight in iterated commutators.  (iii) follows from substitution into~\eqref{eq:thqg-I-hs-criterion-bound}.
\end{proof}

\subsubsection{Convergence radii and analytic continuation}

\begin{proposition}[Convergence radii of the scalar partition function; \ClaimStatusProvedHere]
% label removed: prop:thqg-I-convergence-radii
\index{convergence radius!scalar partition function}
The scalar gravitational partition function $Z^{\mathrm{scal}}_{\mathrm{grav}}(\cA;\,\hbar) = \kappa \cdot \frac{\sqrt{\hbar}/2}{\sin(\sqrt{\hbar}/2)}$ has:
\begin{enumerate}[label=\textup{(\roman*)}]
\item \emph{Radius of convergence in $\hbar$:} $R = 4\pi^2 \approx 39.48$.
\item \emph{First singularity:} A pole of order $1$ at $\hbar = 4\pi^2$, corresponding to the first zero of $\sin(\sqrt{\hbar}/2)$ at $\sqrt{\hbar} = 2\pi$.
\item \emph{Analytic continuation:} $Z^{\mathrm{scal}}$ extends to a meromorphic function on $\mathbb{C}$ with poles at $\hbar = 4\pi^2 n^2$, $n \in \mathbb{Z}_{> 0}$.
\end{enumerate}
\end{proposition}

\begin{proof}
The function $f(\hbar) = \frac{\sqrt{\hbar}/2}{\sin(\sqrt{\hbar}/2)}$ is analytic at $\hbar = 0$ (removable singularity: $f(0) = 1$).  The root test gives $\limsup_{g \to \infty} |a_g|^{1/g} = 1/(2\pi)^2 = 1/(4\pi^2)$, so the radius of convergence is $4\pi^2$.  The poles at $\hbar = 4\pi^2 n^2$ are simple with residues $(-1)^{n+1} \cdot 2n\pi \cdot \kappa$.
\end{proof}

\begin{remark}[Hagedorn singularities]
% label removed: rem:thqg-I-hagedorn
\index{Hagedorn singularity}
The poles at $\hbar = 4\pi^2 n^2$ are the gravitational analogues of the Hagedorn temperature in string theory.  At $\hbar = 4\pi^2$, the genus sum first diverges.  Beyond this point, the partition function requires analytic continuation and the genus expansion breaks down.  In the holographic context, the pole at $\hbar = 4\pi^2$ corresponds to a phase transition in the gravitational theory: the genus expansion becomes unreliable, and non-perturbative effects (saddle-point contributions from higher-topology geometries) must be included.
\end{remark}

\subsubsection{Admissible levels and singular fibers}

\begin{remark}[Admissible levels]
% label removed: rem:thqg-I-admissible
\index{admissible level!finiteness}
At admissible levels $k = -h^\vee + p/q$ (with $p,q$ coprime positive integers and $p \geq h^\vee$), the universal vertex algebra $V_k(\fg)$ has null vectors and the simple quotient $L_k(\fg)$ differs from $V_k(\fg)$.  The HS-sewing criterion applies to $V_k(\fg)$ unconditionally (by Proposition~\ref{prop:thqg-I-standard-hs}).  For the simple quotient $L_k(\fg)$, the growth remains subexponential by the Kac--Wakimoto character formula~\cite{KW88}, so HS-sewing holds as well.

The modular characteristic is $\kappa(V_k(\fg)) = td/(2h^\vee)$ with $t = k + h^\vee$, which vanishes at critical level $k = -h^\vee$.  At critical level, the shadow free energies all vanish: $F_g = 0$ for all $g \geq 1$.  The Feigin--Frenkel center emerges as the replacement for the genus tower.
\end{remark}

\begin{remark}[Singular fibers]
% label removed: rem:thqg-I-singular-fibers
\index{singular fiber!modular characteristic}
As the level $k$ varies, the modular characteristic $\kappa(k) = td/(2h^\vee)$ with $t = k + h^\vee$ traces a linear function of $k$.  The critical level $k = -h^\vee$ is the unique singular fiber where $\kappa$ vanishes.  Near the singular fiber, $\kappa \to 0$ and the shadow free energies $F_g \to 0$ uniformly: $|F_g(k)| \leq C |k + h^\vee| / (2\pi)^{2g}$ for all $g$.  The shadow tower collapses to the trivial element $\Theta_\cA = 0$ at $k = -h^\vee$.
\end{remark}

\begin{remark}[Critical level and the Feigin--Frenkel center]
% label removed: rem:thqg-I-ff-center
\index{Feigin--Frenkel center}
At critical level $k = -h^\vee$, the Sugawara construction is \emph{undefined} (not ``$c$ diverges''; the Sugawara tensor $T(z) = \frac{1}{2(k+h^\vee)} \sum_a :J^a(z) J^a(z):$ has a pole at $k = -h^\vee$).  The Feigin--Frenkel theorem identifies the center $Z(V_{-h^\vee}(\fg)) \cong \mathrm{Fun}\,\mathrm{Op}_{\fg^\vee}(\mathbb{D})$ with the algebra of functions on the space of $\fg^\vee$-opers on the formal disc.  The shadow free energy vanishes at this level because $\kappa = 0$, but the center provides a replacement: the ``classical'' genus-$g$ invariant is the space of $\fg^\vee$-opers on $\Sigma_g$, which is finite-dimensional and does not require the HS-sewing machinery.
\end{remark}

\subsubsection{Summary table}

\begin{table}[ht]
\centering
\caption{Perturbative finiteness: standard landscape}
% label removed: tab:thqg-I-finiteness-summary
\index{perturbative finiteness!summary table}
\renewcommand{\arraystretch}{1.4}
{\footnotesize
\begin{tabular}{|l|c|c|c|c|c|}
\hline
\textbf{Family $\cA$} & $\boldsymbol{\dim H_n}$ & \textbf{OPE $N$} & $\boldsymbol{\kappa}$ & $\boldsymbol{r_{\max}}$ & \textbf{HS Status} \\
\hline
$\mathcal{H}_k$ & $p(n)$ & $1$ & $k$ & $2$ & \textbf{Proved} \\
\hline
$\widehat{\fg}_k$ & $\sim e^{\pi\sqrt{2dn/3}}$ & $1$ & $\frac{(k+h^\vee)d}{2h^\vee}$ & $3$ & \textbf{Proved} \\
\hline
$\mathrm{Vir}_c$ & $p(n)$ & $2$ & $c/2$ & $\infty$ & \textbf{Proved} \\
\hline
$\beta\gamma$ & $p(n)^2$ & $1$ & $1$ & $4$ & \textbf{Proved} \\
\hline
$\mathcal{W}_N$ & $N$-colored & $2N$ & $c_N/2 \cdot \rho_N$ & $\infty$ & \textbf{Proved} \\
\hline
$V_\Lambda$ & $|\Lambda^*/\Lambda|\cdot p_r(n)$ & $r$ & $\rank(\Lambda)$ & $2$ & \textbf{Proved} \\
\hline
Free fermions & polynomial & $0$ & $n/4$ & $2$ & \textbf{Proved} \\
\hline
\end{tabular}}
\end{table}

\begin{remark}[Exhaustiveness]
% label removed: rem:thqg-I-exhaustiveness
\index{standard landscape!exhaustiveness}
The seven families in Table~\ref{tab:thqg-I-finiteness-summary} exhaust the ``standard landscape'' of finitely strongly generated chiral algebras with simple Lie symmetry.  Finiteness extends to:
\begin{enumerate}[label=\textup{(\alph*)}]
\item \emph{Tensor products:} $\cA \otimes \cB$ satisfies HS-sewing if both $\cA$ and $\cB$ do (the HS norm of the tensor product is the product of HS norms).
\item \emph{Cosets:} $\cA / \cB$ (the coset of $\cB \subset \cA$) satisfies HS-sewing if $\cA$ does (the sector dimensions can only decrease).
\item \emph{Orbifolds:} $\cA^G$ (orbifold by a finite group $G$) satisfies HS-sewing if $\cA$ does (the twisted sectors contribute additional terms, but all with subexponential growth).
\item \emph{Non-principal $\mathcal{W}$-algebras:} $\mathcal{W}_k(\fg, f)$ for nilpotent $f$ satisfies HS-sewing at generic $k$ (the DS reduction does not increase OPE degree).
\end{enumerate}
\end{remark}


% ======================================================================
%  SUBSECTION 6: IMPLICATIONS FOR 3D QUANTUM GRAVITY
% ======================================================================

\subsection{Implications for 3d quantum gravity}
% label removed: subsec:thqg-I-3d-gravity
\index{3d quantum gravity|textbf}
\index{BTZ black hole!partition function}

Three-dimensional gravity with negative cosmological constant $\Lambda = -1/\ell^2$ is the physical arena where the perturbative finiteness theorem has the most direct consequences.  The boundary theory is a 2d CFT with central charge determined by the Brown--Henneaux formula $c = 3\ell/(2G)$.

\subsubsection{The Brown--Henneaux identification}

\begin{setup}[Brown--Henneaux central charge]
% label removed: setup:thqg-I-brown-henneaux
\index{Brown--Henneaux!central charge}
The asymptotic symmetry algebra of AdS$_3$ gravity at spatial infinity is two copies of the Virasoro algebra with central charge
\begin{equation}% label removed: eq:thqg-I-brown-henneaux
c \;=\; \frac{3\ell}{2G}\,,
\end{equation}
where $\ell$ is the AdS radius and $G$ is the three-dimensional Newton constant~\cite{BH86}.  In the twisted holographic setting, $\cA = \mathrm{Vir}_c$ is the boundary chiral algebra, with modular characteristic $\kappa(\cA) = c/2 = 3\ell/(4G)$.
\end{setup}

\begin{remark}[Semiclassical limit]
% label removed: rem:thqg-I-semiclassical
\index{semiclassical limit}
The semiclassical limit is $G \to 0$ with $\ell$ fixed, which gives $c \to \infty$.  In this limit, $\kappa \to \infty$ and the shadow free energies $F_g = \kappa \cdot \lambda_g^{\mathrm{FP}} \to \infty$ for each $g \geq 1$: the genus corrections are large in absolute terms but small \emph{relative} to the tree-level contribution $F_0 = \kappa$.  The ratio $F_g/F_0 = \lambda_g^{\mathrm{FP}}$ is independent of $\kappa$ and decreases super-exponentially in $g$.  The genus expansion is therefore an asymptotic expansion in $1/\kappa \sim G/\ell$ that happens to converge (as a formal power series in $\hbar$).
\end{remark}

\subsubsection{The BTZ partition function from shadow free energies}

\begin{theorem}[BTZ partition function from shadow theory; \ClaimStatusProvedHere]
% label removed: thm:thqg-I-btz
\index{BTZ black hole!shadow free energy}
The one-loop correction to the BTZ partition function in the twisted holographic framework is
\begin{equation}% label removed: eq:thqg-I-btz-F1
F_1^{\mathrm{BTZ}} \;=\; F_1(\mathrm{Vir}_c) \;=\; \frac{c}{48} \;=\; \frac{\ell}{32G}\,.
\end{equation}
The full perturbative BTZ partition function, including all genus corrections, is
\begin{equation}% label removed: eq:thqg-I-btz-full
Z_{\mathrm{BTZ}}^{\mathrm{pert}}(\hbar) \;=\; \frac{c}{2} \cdot \frac{\sqrt{\hbar}/2}{\sin(\sqrt{\hbar}/2)} \;=\; \frac{3\ell}{4G} \cdot \frac{\sqrt{\hbar}/2}{\sin(\sqrt{\hbar}/2)}\,.
\end{equation}
\end{theorem}

\begin{proof}
By the Brown--Henneaux identification (Setup~\ref{setup:thqg-I-brown-henneaux}), the boundary algebra is $\mathrm{Vir}_c$ with $\kappa = c/2$.  The genus-$1$ free energy is $F_1 = \kappa/24 = c/48$.  The full perturbative series is the generating function~\eqref{eq:thqg-I-gf-ahat} with $\kappa = c/2$.  Substituting $c = 3\ell/(2G)$ gives $\kappa = 3\ell/(4G)$.
\end{proof}

\begin{remark}[Comparison with known results]
% label removed: rem:thqg-I-comparison-btz
\index{BTZ black hole!comparison}
The genus-$1$ correction $F_1 = c/48$ has been computed by multiple methods: (a)~gravity path integral (Giombi--Maloney--Yin~\cite{GMY08}): the one-loop determinant of the graviton on the solid torus gives $\prod_{n=2}^{\infty} (1 - q^n)^{-1}$; (b)~Chern--Simons theory (Witten~\cite{Witten88}): 3d gravity as $\mathrm{SL}(2,\mathbb{R}) \times \mathrm{SL}(2,\mathbb{R})$ CS theory, with $c/24$ for the full theory and $c/48$ for the holomorphic half; (c)~modular bootstrap (Maloney--Witten~\cite{MW10}).  All agree with the shadow free energy computation.  The advantage of the shadow method is that it extends automatically to all genera via the closed form~\eqref{eq:thqg-I-btz-full}.
\end{remark}

\subsubsection{Higher-genus corrections}

\begin{proposition}[Higher-genus BTZ corrections; \ClaimStatusProvedHere]
% label removed: prop:thqg-I-btz-higher-genus
\index{BTZ black hole!higher-genus corrections}
The genus-$g$ correction to the BTZ partition function in the scalar channel is
\begin{equation}% label removed: eq:thqg-I-btz-Fg
F_g^{\mathrm{BTZ}} \;=\; \frac{c}{2} \cdot \frac{2^{2g-1}-1}{2^{2g-1}} \cdot \frac{|B_{2g}|}{(2g)!} \;=\; \frac{3\ell}{4G} \cdot \lambda_g^{\mathrm{FP}}\,.
\end{equation}
The first five corrections are:
\begin{align}
F_1^{\mathrm{BTZ}} &= \frac{c}{48}\,, % label removed: eq:thqg-I-btz-F1-val \\
F_2^{\mathrm{BTZ}} &= \frac{7c}{11520}\,, % label removed: eq:thqg-I-btz-F2-val \\
F_3^{\mathrm{BTZ}} &= \frac{31c}{1935360}\,, % label removed: eq:thqg-I-btz-F3-val \\
F_4^{\mathrm{BTZ}} &= \frac{127c}{309657600}\,, % label removed: eq:thqg-I-btz-F4-val \\
F_5^{\mathrm{BTZ}} &= \frac{73c}{7007109120}\,. % label removed: eq:thqg-I-btz-F5-val
\end{align}
The ratio $F_{g+1}/F_g \to 1/(4\pi^2)$ as $g \to \infty$, consistent with the Hagedorn singularity at $\hbar = 4\pi^2$.
\end{proposition}

\begin{proof}
Direct from $\kappa(\mathrm{Vir}_c) = c/2$ and the universal formula (Proposition~\ref{prop:thqg-I-Fg-closed-form}).
\end{proof}

\begin{remark}[Numerical values for $c = 24$ (monster CFT)]
% label removed: rem:thqg-I-monster
\index{monster CFT}
For the monster CFT with $c = 24$ (the Frenkel--Lepowsky--Meurman moonshine module $V^\natural$), $\kappa = 12$:
\begin{align*}
F_1 &= 12/24 = 1/2, \\
F_2 &= 7 \cdot 12/5760 = 84/5760 = 7/480, \\
F_3 &= 31 \cdot 12/967680 = 372/967680 = 31/80640, \\
F_4 &= 127 \cdot 12/154828800 = 1524/154828800 = 127/12902400.
\end{align*}
The one-loop contribution $F_1 = 1/2$ is order $1$, and each subsequent genus is suppressed by the Bernoulli factor.  The total perturbative partition function at $\hbar = 1$ is $Z = 12 \cdot 1/\sin(1/2) \approx 12 \cdot 2.086 \approx 25.0$.
\end{remark}

\subsubsection{The quartic contact invariant in sub-subleading corrections}

\begin{remark}[Quartic contact invariant in 3d gravity]
% label removed: rem:thqg-I-quartic-3d
\index{quartic contact invariant!3d gravity}
The scalar channel captures only the arity-$0$ information.  The first non-scalar correction comes from the quartic contact invariant
\begin{equation}% label removed: eq:thqg-I-quartic-contact
Q^{\mathrm{contact}}_{\mathrm{Vir}} \;=\; \frac{10}{c(5c+22)}\,,
\end{equation}
which is the arity-$4$, genus-$0$ component of the shadow tower.  For $c = 3\ell/(2G)$, this gives $Q^{\mathrm{contact}} = 10/(c(5c+22)) = 2/c^2 + O(1/c^3)$, a two-loop correction in Newton's constant.  The genus-$1$ Hessian correction is
\begin{equation}% label removed: eq:thqg-I-genus1-hessian
\delta_H^{(1)}_{\mathrm{Vir}} \;=\; \frac{120}{c^2(5c+22)}\,x^2\,.
\end{equation}
These corrections are specific to the mixed shadow class ($r_{\max} = \infty$) of the Virasoro algebra.
\end{remark}

\begin{remark}[Shadow depth and gravitational complexity]
% label removed: rem:thqg-I-depth-complexity
\index{shadow depth!gravitational complexity}
The shadow depth classification ($G/L/C/M$ for Gaussian/Lie/Contact/Mixed) stratifies 3d gravity theories by their nonlinear complexity:
\begin{enumerate}[label=\textup{(\alph*)}]
\item \emph{Gaussian ($r_{\max} = 2$):} Heisenberg, lattice VOAs, free fermions.  The gravitational theory is free.  The partition function is an exact Fredholm determinant.  No nonlinear corrections.
\item \emph{Lie-type ($r_{\max} = 3$):} Affine Kac--Moody algebras.  Cubic interactions (tree-level scattering).  The partition function receives corrections from trivalent graphs.
\item \emph{Contact-type ($r_{\max} = 4$):} $\beta\gamma$ system.  Quartic interactions but no higher.  The shadow tower terminates after two nonlinear steps.
\item \emph{Mixed ($r_{\max} = \infty$):} Virasoro, $\mathcal{W}_N$.  Interactions at all orders.  The full shadow tower is infinite, and the perturbative finiteness theorem is needed in full strength.  The quintic obstruction $o^{(5)}_{\mathrm{Vir}} \neq 0$ forces the tower to be genuinely infinite.
\end{enumerate}
The Brown--Henneaux boundary theory for pure 3d gravity is $\mathrm{Vir}_c$, which falls in the mixed class.
\end{remark}

\subsubsection{Large-$c$ expansion and gravitational coupling}

\begin{proposition}[Leading gravitational coupling expansion; \ClaimStatusProvedHere]
% label removed: prop:thqg-I-1c-expansion
\index{gravitational coupling!large-$c$ expansion}
In the semiclassical expansion $\hbar = 1/c$, the perturbative gravitational partition function has the asymptotic form
\begin{equation}% label removed: eq:thqg-I-1c-expansion
\log Z_{\mathrm{grav}}^{\mathrm{scal}}(\mathrm{Vir}_c;\,1/c) \;=\; \frac{1}{48} + \frac{7}{11520\,c} + \frac{31}{1935360\,c^2} + \frac{127}{309657600\,c^3} + O(c^{-4})\,.
\end{equation}
The leading term $1/48$ is the one-loop correction, and each subsequent term is a higher-loop gravitational correction suppressed by one additional power of $G/\ell \sim 1/c$.
\end{proposition}

\begin{proof}
From $F_g(\mathrm{Vir}_c) = (c/2)\lambda_g^{\mathrm{FP}}$, the contribution of genus~$g$ to $\log Z$ at $\hbar = 1/c$ is $\lambda_g^{\mathrm{FP}}/(2c^{g-1})$.  The expansion $\log Z = \sum_{g \geq 1} F_g \hbar^g = \sum_{g \geq 1} (c/2)\lambda_g^{\mathrm{FP}}/c^g = \sum_{g \geq 1} \lambda_g^{\mathrm{FP}} / (2c^{g-1})$.  At $g = 1$: $\lambda_1^{\mathrm{FP}} = 1/24$, giving $1/(2 \cdot 24) = 1/48$.  At $g = 2$: $7/(2 \cdot 5760 \cdot c) = 7/(11520c)$.  And so on.
\end{proof}

\begin{proposition}[Newton's constant expansion; \ClaimStatusProvedHere]
% label removed: prop:thqg-I-newton-expansion
\index{Newton's constant!expansion}
Substituting $c = 3\ell/(2G)$ and $\hbar = 2G/(3\ell)$, the gravitational partition function becomes
\begin{equation}% label removed: eq:thqg-I-newton-expansion
Z_{\mathrm{grav}}^{\mathrm{scal}}(\mathrm{Vir}_c;\,2G/(3\ell)) \;=\; \frac{3\ell}{4G} \cdot \frac{\sqrt{2G/(3\ell)}/2}{\sin(\sqrt{2G/(3\ell)}/2)}\,.
\end{equation}
For $G/\ell \ll 1$ (weak gravitational coupling), the expansion in $G/\ell$ gives
\begin{equation}% label removed: eq:thqg-I-G-expansion
Z \;=\; \frac{3\ell}{4G}\left(1 + \frac{G}{72\ell} + \frac{7G^2}{51840\ell^2} + O(G^3/\ell^3)\right).
\end{equation}
\end{proposition}

\begin{proof}
Direct Taylor expansion of $x/(2\sin(x/2))$ at $x = 0$ with $x^2 = 2G/(3\ell)$.
\end{proof}

\subsubsection{Towards the non-perturbative regime}

\begin{conjecture}[Non-perturbative completion; \ClaimStatusConjectured]
% label removed: conj:thqg-I-non-perturbative
\index{non-perturbative!gravitational partition function}
The meromorphic function $Z^{\mathrm{scal}}_{\mathrm{grav}}(\hbar) = \kappa \cdot \frac{\sqrt{\hbar}/2}{\sin(\sqrt{\hbar}/2)}$ admits a non-perturbative completion satisfying:
\begin{enumerate}[label=\textup{(\roman*)}]
\item Resurgent structure: the Borel transform has singularities at $\hbar = 4\pi^2 n^2$.
\item Stokes phenomena: crossing Stokes rays changes the partition function by $\sim e^{-4\pi^2 n^2/\hbar}$, identified with non-handlebody saddle points.
\item Positivity: for real $\hbar > 0$, the full non-perturbative partition function $Z > 0$.
\end{enumerate}
\end{conjecture}

\begin{evidence}[Evidence for the resurgent structure]
% label removed: evid:thqg-I-resurgent
The Taylor coefficients $F_g \sim 2\kappa/(2\pi)^{2g}$ have the form $a_g \sim C \cdot A^{-g}$ with $A = (2\pi)^2$, which is precisely the asymptotic form associated with a simple pole in the Borel plane at $\xi = 4\pi^2$.  The Borel transform is
\begin{equation}% label removed: eq:thqg-I-borel-transform
\hat{Z}(\xi) \;=\; \sum_{g=1}^\infty \frac{F_g}{\Gamma(g)} \xi^{g-1} \;=\; \kappa \sum_{g=1}^\infty \frac{\lambda_g^{\mathrm{FP}}}{(g-1)!} \xi^{g-1}.
\end{equation}
Since $\lambda_g^{\mathrm{FP}} \sim 2/(2\pi)^{2g}$, the ratio $\lambda_g^{\mathrm{FP}}/\Gamma(g) \sim 2/((2\pi)^{2g}(g-1)!)$, and the series converges for all $\xi \in \mathbb{C}$ (the Borel transform is entire).  The singularities at $\xi = 4\pi^2 n^2$ in the Borel plane arise from the Laplace integral relating the original series to the Borel sum: the non-perturbative corrections $\sim e^{-4\pi^2 n^2/\hbar}$ match the mass gap of the lightest BTZ black hole above the vacuum.
\end{evidence}

\begin{remark}[BTZ black holes as non-perturbative saddles]
% label removed: rem:thqg-I-btz-saddles
\index{BTZ black hole!non-perturbative}
In 3d gravity, the BTZ black holes of mass $M = n^2/(8G\ell)$ ($n \geq 1$) are the non-perturbative saddle points.  Their Euclidean action is $I_n = 4\pi^2 n^2 / \hbar$ (with $\hbar = 2G/(3\ell)$).  The Stokes phenomenon at the $n$-th pole $\hbar = 4\pi^2 n^2$ corresponds to the creation/annihilation of the $n$-th BTZ black hole as one crosses the Stokes ray.  The full non-perturbative partition function should be $Z = Z_{\mathrm{pert}} + \sum_{n \geq 1} c_n e^{-I_n/\hbar}$ with explicit coefficients $c_n$ determined by the resurgent structure.
\end{remark}

%% ---- Stokes phenomena on the Steinberg object ----

\providecommand{\SteinbergB}{\mathfrak{S}_b}

\subsubsection{Stokes phenomena on the Steinberg object}
% label removed: sssec:thqg-stokes-steinberg
\index{Stokes phenomena!Steinberg object}
\index{Steinberg object!Stokes filtration}

The poles of $Z_{\mathrm{grav}}(\hbar) = \kappa \cdot \frac{\sqrt{\hbar}/2}{\sin(\sqrt{\hbar}/2)}$ at $\hbar = 4\pi^2 n^2$ are not merely analytic accidents: they are the Stokes lines in the Borel plane of the convolution integral on the Steinberg object~$\SteinbergB$.  The Borel transform
\[
B[Z](\zeta) \;=\; \sum_{g \geq 1} \frac{F_g}{\Gamma(g)}\,\zeta^{g-1}
\]
has singularities at $\zeta_n = 4\pi^2 n^2$, which are the Euclidean actions of the $n$-th BTZ saddle points (Remark~\ref{rem:thqg-I-btz-saddles}).  The Stokes-filtered Steinberg construction upgrades the perturbative Steinberg presentation to one that sees all saddle-point sectors simultaneously.

\begin{construction}[The Stokes-filtered Steinberg]
% label removed: constr:thqg-stokes-filtered-steinberg
\index{Steinberg object!Stokes-filtered}
Define the \emph{Stokes-filtered Steinberg} $\SteinbergB^{\mathrm{Stokes}}$ as follows.  For each $n \geq 1$, let $\SteinbergB^{\mathrm{BTZ},n}$ denote the Steinberg object associated to the $n$-th BTZ saddle, i.e.\ the analogue of~$\SteinbergB$ constructed from the gauge equivalence data of the $n$-th non-handlebody saddle point with Euclidean action $I_n = 4\pi^2 n^2 / \hbar$.  The perturbative Steinberg~$\SteinbergB$ is $\SteinbergB^{\mathrm{BTZ},0}$ (the vacuum saddle).  The Stokes-filtered Steinberg is the datum
\[
\SteinbergB^{\mathrm{Stokes}} \;=\; \bigl(\SteinbergB,\;\{\SteinbergB^{\mathrm{BTZ},n}\}_{n \geq 1},\;\{S_n\}_{n \geq 1}\bigr)
\]
where $S_n \colon H_*^{\mathrm{BM}}(\SteinbergB) \to H_*^{\mathrm{BM}}(\SteinbergB^{\mathrm{BTZ},n})$ is the Stokes map induced by analytic continuation of the Borel-resummed integral across the $n$-th Stokes ray $\arg(\zeta) = \arg(\zeta_n)$.
\end{construction}

\begin{proposition}[Resurgent structure from Stokes-filtered BM homology; \ClaimStatusConjectured]
% label removed: prop:thqg-stokes-bm-resurgence
\index{resurgence!Steinberg BM homology}
\index{BM homology!Stokes-filtered}
The resurgent structure of $Z_{\mathrm{grav}}$ is presented by the Stokes-filtered Borel--Moore homology of the Steinberg:
\begin{equation}% label removed: eq:thqg-stokes-bm-decomposition
H_*^{\mathrm{BM},\,\mathrm{Stokes}}(\SteinbergB) \;=\; H_*^{\mathrm{BM}}(\SteinbergB) \;\oplus\; \bigoplus_{n \geq 1} H_*^{\mathrm{BM}}(\SteinbergB^{\mathrm{BTZ},n})\,,
\end{equation}
where $\SteinbergB^{\mathrm{BTZ},n}$ is the Steinberg object for the $n$-th BTZ saddle.  The Stokes automorphism
\[
\mathfrak{S}_n \colon H_*^{\mathrm{BM}}(\SteinbergB) \;\longrightarrow\; H_*^{\mathrm{BM}}(\SteinbergB)
\]
is the wall-crossing automorphism exchanging the perturbative and $n$-th BTZ sectors.  More precisely, the discontinuity of the lateral Borel resummation across the $n$-th Stokes ray is
\[
\mathrm{disc}_n\, Z \;=\; \mathfrak{S}_n([\SteinbergB]) \cdot e^{-\zeta_n / \hbar}\,,
\]
where $[\SteinbergB] \in H_*^{\mathrm{BM}}(\SteinbergB)$ is the fundamental class and $\zeta_n = 4\pi^2 n^2$.
\end{proposition}

\begin{remark}[Positivity from geometry]
% label removed: rem:thqg-stokes-positivity
\index{positivity!gravitational partition function}
\index{BM homology!fundamental class positivity}
The positivity of $Z_{\mathrm{grav}}$ for real $\hbar > 0$ (away from the poles $\hbar = 4\pi^2 n^2$) admits a geometric explanation through the Steinberg presentation.  The perturbative contribution is $\langle [\SteinbergB], [\SteinbergB] \rangle_{\mathrm{BM}}$, the self-pairing of the BM fundamental class $[\SteinbergB]$, which is positive since $\SteinbergB$ carries a canonical orientation as a complex-algebraic variety.  The Stokes phenomena at $\hbar = 4\pi^2 n^2$ are phase transitions where the $n$-th BTZ stratum $\SteinbergB^{\mathrm{BTZ},n}$ becomes visible: as $\hbar$ crosses $4\pi^2 n^2$, the Borel resummation contour passes through the singularity $\zeta_n$, and the new Steinberg stratum contributes an additional sector to the partition function.  The pole itself reflects the divergence of the lateral resummation integral at the singular point: geometrically, the collision of the perturbative fundamental cycle with the $n$-th BTZ cycle in the Stokes-filtered homology.
\end{remark}

\begin{remark}[Pure 3d gravity as a model]
% label removed: rem:thqg-I-pure-3d
\index{3d gravity!pure}
Pure 3d gravity with $\Lambda < 0$ is the simplest non-trivial gravitational theory.  Its boundary algebra is $\mathrm{Vir}_c$ with $c = 3\ell/(2G)$, which falls in the mixed class.  The perturbative finiteness theorem provides a complete, explicit, genus-by-genus calculation of the gravitational partition function:
\[
Z_{\mathrm{grav}}^{\mathrm{scal}}(\hbar) \;=\; \frac{3\ell}{4G} \cdot \frac{\sqrt{\hbar}/2}{\sin(\sqrt{\hbar}/2)}\,.
\]
The non-scalar corrections from $Q^{\mathrm{contact}} = 10/(c(5c+22))$ and the infinite shadow tower encode the genuinely nonlinear aspects of the gravitational theory, suppressed by $1/c^2$ relative to the scalar contribution.
\end{remark}

\begin{remark}[Comparison with matrix models]
% label removed: rem:thqg-I-matrix-models
\index{matrix model!gravitational partition function}
In 2d topological gravity (Witten--Kontsevich), the genus-$g$ free energy involves Bernoulli numbers with a different prefactor: the Witten--Kontsevich generating function is $\sum_g F_g^{\mathrm{WK}} (2g-2)! x^{2g-2}$, with $F_g^{\mathrm{WK}} = \chi(\mathcal{M}_g) / (2g-2)!$ involving the orbifold Euler characteristic.  The relation to the shadow generating function is $\frac{x/2}{\sin(x/2)} = \left.\frac{x/2}{\sinh(x/2)}\right|_{x \mapsto ix}$, reflecting the Wick rotation from Lorentzian to Euclidean signature.  In the matrix model, the $\sinh$ function arises from the Gaussian integral, while in the shadow theory, the $\sin$ function arises from the Hodge integral; the two are related by analytic continuation $x \mapsto ix$.
\end{remark}

\begin{remark}[Higher-spin gravity]
% label removed: rem:thqg-I-higher-spin
\index{higher-spin gravity}
For higher-spin gravity in AdS$_3$ with $\mathcal{W}_N$ symmetry, the boundary algebra is $\mathcal{W}_k(\mathfrak{sl}_N)$ at $c = (N-1)(1 - N(N+1)/(k+N))$.  The modular characteristic is $\kappa(\mathcal{W}_N) = c/2 \cdot \rho_N$ where $\rho_N$ is a normalizing factor.  The perturbative partition function is
\[
Z_{\mathrm{grav}}^{W_N}(\hbar) \;=\; \kappa(\mathcal{W}_N) \cdot \frac{\sqrt{\hbar}/2}{\sin(\sqrt{\hbar}/2)}\,,
\]
which has the same functional form as the Virasoro case but with a different normalization.  The qualitative features (convergence for $|\hbar| < 4\pi^2$, Hagedorn singularity, resurgent structure) are identical.  The difference lies in the non-scalar corrections: $\mathcal{W}_N$ has shadow depth $r_{\max} = \infty$ and the shadow tower includes contributions from all $W^{(s)}$ fields ($s = 2, \ldots, N$).
\end{remark}

\subsubsection{The holographic dictionary}

\begin{remark}[The holographic dictionary]
% label removed: rem:thqg-I-holographic-dictionary
\index{holographic dictionary!perturbative finiteness}
The perturbative finiteness theorem (Theorem~\ref{thm:thqg-I-perturbative-finiteness}) translates into the holographic dictionary:
\begin{center}
\renewcommand{\arraystretch}{1.3}
\begin{tabular}{l|l}
\textbf{Boundary (algebra)} & \textbf{Bulk (gravity)} \\
\hline
$\cA$ (chiral algebra) & boundary conditions \\
$\kappa(\cA)$ (modular characteristic) & cosmological constant \\
$F_g(\cA)$ (shadow free energy) & genus-$g$ path integral \\
HS-sewing condition & UV finiteness \\
$\Theta_\cA$ (MC element) & classical solution + perturbations \\
$D_\cA^2 = 0$ & Ward identity \\
$r_{\max}(\cA)$ (shadow depth) & gravitational complexity \\
$4\pi^2$ (convergence radius) & Hagedorn temperature \\
$Q^{\mathrm{contact}}$ (quartic invariant) & two-loop graviton self-energy \\
$\cA^!$ (Koszul dual) & dual boundary conditions \\
\end{tabular}
\end{center}
The passage from boundary data to bulk partition function is algorithmic: given the OPE data of $\cA$, the modular characteristic $\kappa$ is computable, the Faber--Pandharipande coefficients $\lambda_g^{\mathrm{FP}}$ are universal, and the partition function $Z_{\mathrm{grav}} = \kappa \cdot x/(2\sin(x/2))$ follows.  No UV regularization, no renormalization group flow, and no anomaly matching are required.
\end{remark}

\begin{remark}[Summary of the finiteness mechanism]
% label removed: rem:thqg-I-finiteness-summary
\index{perturbative finiteness!mechanism summary}
The perturbative finiteness of twisted quantum gravity rests on a chain of three implications:
\begin{equation}% label removed: eq:thqg-I-chain
\underbrace{D_\cA^2 = 0}_{\text{bar complex}} \;\Longrightarrow\; \underbrace{\Theta_\cA \in \mathrm{MC}(\gAmod)}_{\text{exact MC}} \;\Longrightarrow\; \underbrace{F_g = \kappa \cdot \lambda_g^{\mathrm{FP}} < \infty}_{\text{finite amplitudes}}.
\end{equation}
The first implication is the bar-intrinsic construction (Theorem~\ref*{V1-thm:mc2-bar-intrinsic}): the element $\Theta_\cA := D_\cA - d_0$ is automatically Maurer--Cartan because the bar differential squares to zero.  The second implication is the projection onto genus~$g$: the shadow free energy is the integral of a cohomology class over a compact moduli space, hence finite.  The entire mechanism is algebraic (no analysis of divergences, no regularization, no renormalization) and exact (no perturbative truncation, no asymptotic expansion).
\end{remark}


% ======================================================================
%  SUBSECTION 7: THE HS-SEWING CONDITION AND THE MC RECURSION
% ======================================================================

\subsection{HS-sewing and the MC recursion}
% label removed: subsec:thqg-I-hs-mc-recursion
\index{HS-sewing!MC recursion}

The HS-sewing condition and the MC recursion interact at two levels: the HS condition provides the \emph{analytic} input (convergence of sewing amplitudes), while the MC recursion provides the \emph{algebraic} input (unique determination of higher-genus data from lower-genus data).  Together, they give a complete genus-by-genus algorithm for computing the gravitational partition function.

\subsubsection{The sewing amplitude as a graph sum}

\begin{proposition}[Graph-sum representation of sewing amplitudes; \ClaimStatusProvedHere]
% label removed: prop:thqg-I-graph-sum-sewing
\index{sewing amplitude!graph-sum representation}
Let $\cA$ be a modular Koszul chiral algebra satisfying HS-sewing.  The genus-$g$ sewing amplitude on a surface $\Sigma_g$ with pants decomposition $\mathcal{P}$ has the graph-sum representation
\begin{equation}% label removed: eq:thqg-I-graph-sewing
Z_{\mathcal{P}}(\cA;\,\mathbf{q}) \;=\; \sum_{\Gamma \in \Gmod(g)} \frac{1}{|\Aut(\Gamma)|}\,\ell_\Gamma(\mathbf{q}),
\end{equation}
where the sum runs over connected modular graphs $\Gamma$ of total genus $g$, and $\ell_\Gamma(\mathbf{q})$ is the graph amplitude: the product of vertex labels (transferred cyclic operations $\ell_n^{\mathrm{tr}}$) and edge propagators ($P_\cA(q_e) = \sum_{n \geq 0} q_e^n \pi_n$, the weighted projection onto the $n$-th conformal weight).
\end{proposition}

\begin{proof}
The pants decomposition $\mathcal{P}$ determines a dual graph $\Gamma_\mathcal{P}$ with one vertex for each pair of pants and one edge for each sewing circle.  The vertex labels are the three-point functions of $\cA$ (the OPE coefficients $C^c_{a,b}$), and the edge labels are the sewing propagators.  The sewing amplitude is the contraction of all vertex and edge tensors:
\[
Z_{\mathcal{P}} \;=\; \sum_{\text{weight assignments}} \prod_{\text{vertices}} C^{c_v}_{a_v,b_v} \prod_{\text{edges}} q_e^{n_e} \delta_{n_e = c_{v'} = a_{v''}},
\]
where the sum over weight assignments runs over all ways to assign conformal weights to the internal edges.  This is precisely the graph amplitude $\ell_{\Gamma_\mathcal{P}}(\mathbf{q})$.

The independence from the pants decomposition (sewing consistency) means that $Z$ depends only on the genus $g$ and the moduli of $\Sigma_g$, not on the choice of $\mathcal{P}$.  Different pants decompositions give different graphs $\Gamma_\mathcal{P}$, but the sum of all amplitudes is the same.
\end{proof}

\begin{proposition}[Convergence of the graph sum at fixed genus; \ClaimStatusProvedHere]
% label removed: prop:thqg-I-graph-sum-convergence
\index{graph sum!convergence}
Under HS-sewing, the graph sum~\eqref{eq:thqg-I-graph-sewing} converges absolutely at each genus $g$.  The number of contributing graphs is finite at each arity level, and the arity sum converges:
\begin{equation}% label removed: eq:thqg-I-graph-sum-arity
|Z_{\mathcal{P}}(\cA;\,\mathbf{q})| \;\leq\; \sum_{r=0}^{\infty} \sum_{\Gamma \in \Gmod(g,r)} \frac{|\ell_\Gamma(\mathbf{q})|}{|\Aut(\Gamma)|} \;<\; \infty.
\end{equation}
\end{proposition}

\begin{proof}
At each arity $r$, the number of graphs $|\Gmod(g,r)|$ is finite (Proposition~\ref{prop:thqg-I-graph-count}).  Each graph amplitude $|\ell_\Gamma(\mathbf{q})|$ is bounded by the product of vertex norms and propagator norms: $|\ell_\Gamma| \leq \prod_{v} \|\ell_{n_v}^{\mathrm{tr}}\| \cdot \prod_{e} q_e^{n_e}$.  The HS-sewing condition bounds the sum over internal weight assignments.  The sum over $r$ converges because the vertex norms $\|\ell_r^{\mathrm{tr}}\|$ decay (by the shadow tower convergence, Theorem~\ref{thm:recursive-existence}) and the number of graphs grows at most combinatorially while each amplitude decays exponentially in $r$ (from the $q$-weighting).
\end{proof}

\subsubsection{Interplay between analytic and algebraic finiteness}

\begin{remark}[Two finiteness mechanisms]
% label removed: rem:thqg-I-two-mechanisms
\index{finiteness!two mechanisms}
The perturbative finiteness theorem draws on two independent mechanisms:
\begin{enumerate}[label=\textup{(\roman*)}]
\item \emph{Algebraic finiteness} (from the MC equation).  The bar-intrinsic construction $\Theta_\cA = D_\cA - d_0$ gives an exact MC element at all genera.  The genus-$g$ component $\Theta^{(g)}$ is uniquely determined by the lower-genus data through the contracting homotopy.  This determines the \emph{cohomology class} $[\operatorname{Sh}_{g,0}(\Theta_\cA)] \in H^*(\overline{\mathcal{M}}_g)$ without reference to convergence.
\item \emph{Analytic finiteness} (from HS-sewing).  The HS condition guarantees that the sewing construction converges and that the genus-$g$ partition function (as a \emph{number}, not just a cohomology class) is finite.  This is needed for the physical interpretation: the partition function must be a convergent integral, not just a formal expression.
\end{enumerate}
The algebraic mechanism gives existence and uniqueness of the genus-$g$ partition function as a formal quantity.  The analytic mechanism gives convergence and finiteness as a number.  Together, they give a complete, finite, unique answer at every genus.
\end{remark}

\begin{proposition}[Consistency of algebraic and analytic finiteness; \ClaimStatusProvedHere]
% label removed: prop:thqg-I-consistency
\index{finiteness!consistency}
For a modular Koszul chiral algebra $\cA$ satisfying HS-sewing:
\begin{enumerate}[label=\textup{(\roman*)}]
\item The formal shadow free energy $F_g^{\mathrm{formal}}(\cA) := \kappa \cdot \lambda_g^{\mathrm{FP}}$ (computed algebraically from the MC element) agrees with the analytic partition function $F_g^{\mathrm{analytic}}(\cA) := \int_{\overline{\mathcal{M}}_g} Z_{\mathcal{P}}(\cA;\,\mathbf{q})\,\mu_{\mathrm{WP}}$ (computed by sewing and integration):
\begin{equation}% label removed: eq:thqg-I-agreement
F_g^{\mathrm{formal}}(\cA) \;=\; F_g^{\mathrm{analytic}}(\cA).
\end{equation}
\item The analytic continuation of the genus series $\sum_g F_g \hbar^g$ to the meromorphic function $\kappa \cdot \sqrt{\hbar}/(2\sin(\sqrt{\hbar}/2))$ is compatible with the Borel resummation of the sewing amplitudes.
\end{enumerate}
\end{proposition}

\begin{proof}
Part~(i).  The formal shadow free energy is computed from the universal MC element $\Theta_\cA$: the genus-$g$ projection onto arity $0$ gives $F_g = \kappa \cdot \int_{\overline{\mathcal{M}}_g} \lambda_g$.  The analytic partition function is computed by sewing: the genus-$g$ amplitude is the trace of the sewing operator, integrated over $\overline{\mathcal{M}}_g$ with the Weil--Petersson measure.  The agreement follows from the identification of the shadow representation $\operatorname{Sh}_{g,0}(\Theta_\cA)$ with the sewing integrand: both are the genus-$g$ vacuum amplitude of $\cA$, computed by different methods (algebraic graph sum versus analytic trace).

Part~(ii).  The Borel resummation of the sewing amplitudes produces a function of $\hbar$ that agrees with $\kappa \cdot \sqrt{\hbar}/(2\sin(\sqrt{\hbar}/2))$ in the domain of convergence $|\hbar| < 4\pi^2$.  The meromorphic extension is unique by the identity theorem for analytic functions.
\end{proof}

\subsubsection{The full arity-summed partition function}

\begin{theorem}[Full partition function convergence; \ClaimStatusProvedHere]
% label removed: thm:thqg-I-full-convergence
\index{gravitational partition function!full convergence}
For a modular Koszul chiral algebra $\cA$ of shadow depth $r_{\max}$ satisfying HS-sewing, the full gravitational partition function (including all arities) satisfies:
\begin{enumerate}[label=\textup{(\roman*)}]
\item If $r_{\max} < \infty$ (Gaussian, Lie, or contact class), then
\begin{equation}% label removed: eq:thqg-I-full-finite-depth
Z_{\mathrm{grav}}(T;\,\hbar) \;=\; \sum_{g=0}^\infty \hbar^g \sum_{r=0}^{r_{\max}} \int_{\overline{\mathcal{M}}_g} \operatorname{Sh}_{g,r}(\Theta_\cA)
\end{equation}
is a finite sum in arity at each genus.  The genus sum converges absolutely for $|\hbar| < 4\pi^2$.
\item If $r_{\max} = \infty$ (mixed class), then the arity sum converges at each genus (by the shadow tower convergence, Theorem~\ref{thm:recursive-existence}), and the genus sum converges absolutely for $|\hbar| < R_{\mathrm{grav}}$ with $R_{\mathrm{grav}} \geq 4\pi^2$.
\item For any shadow depth, the leading term in the genus expansion is the scalar partition function:
\begin{equation}% label removed: eq:thqg-I-leading-term
Z_{\mathrm{grav}}(T;\,\hbar) \;=\; Z_{\mathrm{grav}}^{\mathrm{scal}}(T;\,\hbar) \;+\; O(\hbar/\kappa^2).
\end{equation}
The non-scalar corrections are suppressed by at least $1/\kappa^2$ relative to the scalar term.
\end{enumerate}
\end{theorem}

\begin{proof}
Part~(i) is immediate from the finite shadow depth: at each genus $g$, only arities $r \leq r_{\max}$ contribute, and each contribution is finite.  The genus sum converges because the scalar term (which dominates) has convergence radius $4\pi^2$.

Part~(ii) follows from the shadow tower convergence and the HS-sewing bound.  At each genus $g$, the arity-$r$ contribution is bounded by $C_r \cdot |\kappa|^{r-2} / (2\pi)^{2g}$ (the factor $|\kappa|^{r-2}$ comes from the vertex labels scaling as $\kappa$, and the extra arity costs additional propagator insertions).  The genus sum of the arity-$r$ term converges for $|\hbar| < 4\pi^2$ independently of $r$, and the arity sum converges by the shadow tower.

Part~(iii) follows from the arity-$r$ scaling: the arity-$r$ contribution to $F_g$ is of order $\kappa \cdot (\kappa^{-1})^{r-2} = \kappa^{3-r}$ for $r \geq 2$.  The leading correction is at $r = 2$ (the spectral discriminant), which is of order $\kappa^1$.  But the spectral discriminant enters at genus $g \geq 1$, so the correction to the genus sum is $O(\kappa \cdot \hbar)$.  The quartic contact invariant enters at $r = 4$ with coefficient $Q^{\mathrm{contact}} = O(\kappa^{-1})$, giving a correction of order $O(\hbar/\kappa)$.  The claim follows.
\end{proof}

\subsubsection{Modular characteristics for all standard families}

\begin{table}[ht]
\centering
\caption{Modular characteristics, shadow depths, and shadow depth classes for the standard landscape.}
% label removed: tab:thqg-I-modular-characteristics
\index{modular characteristic!standard landscape}
\renewcommand{\arraystretch}{1.4}
{\footnotesize
\begin{tabular}{|l|c|c|c|c|c|}
\hline
\textbf{Family $\cA$} & $\boldsymbol{\kappa(\cA)}$ & $\boldsymbol{r_{\max}}$ & \textbf{Class} & $\boldsymbol{F_1(\cA)}$ & $\boldsymbol{F_2(\cA)}$ \\
\hline
$\mathcal{H}_k$ (rank $d$) & $dk$ & $2$ & G & $dk/24$ & $7dk/5760$ \\
\hline
$\widehat{\mathfrak{sl}}_{2,k}$ & $\frac{3(k+2)}{4}$ & $3$ & L & $\frac{k+2}{32}$ & $\frac{7(k+2)}{7680}$ \\
\hline
$\widehat{\mathfrak{sl}}_{N,k}$ & $\frac{(k+N)(N^2-1)}{2N}$ & $3$ & L & $\frac{(k+N)(N^2-1)}{48N}$ & $\frac{7(k+N)(N^2-1)}{11520N}$ \\
\hline
$\widehat{\mathfrak{so}}_{2N,k}$ & $\frac{(k+2N-2)N(2N-1)}{2(2N-2)}$ & $3$ & L & --- & --- \\
\hline
$\mathrm{Vir}_c$ & $c/2$ & $\infty$ & M & $c/48$ & $7c/11520$ \\
\hline
$\beta\gamma$ & $1$ & $4$ & C & $1/24$ & $7/5760$ \\
\hline
$\mathcal{W}_3(k)$ & $\frac{c_3}{2} \rho_3$ & $\infty$ & M & --- & --- \\
\hline
$\mathcal{W}_N(k)$ & $\frac{c_N}{2} \rho_N$ & $\infty$ & M & --- & --- \\
\hline
$V_\Lambda$ & $\rank(\Lambda)$ & $2$ & G & $\rank(\Lambda)/24$ & $7\rank(\Lambda)/5760$ \\
\hline
$V_\Lambda^{N,q}$ (twisted) & $\rank(\Lambda)$ & $2$ & G & $\rank(\Lambda)/24$ & $7\rank(\Lambda)/5760$ \\
\hline
Free fermions (rank $n$) & $n/4$ & $2$ & G & $n/96$ & $7n/23040$ \\
\hline
$\widehat{\mathfrak{sl}}_{2,k}^{\otimes d}$ & $\frac{3d(k+2)}{4}$ & $3$ & L & $\frac{d(k+2)}{32}$ & $\frac{7d(k+2)}{7680}$ \\
\hline
\end{tabular}}
\end{table}

\begin{remark}[Lattice curvature-braiding orthogonality]
% label removed: rem:thqg-I-lattice-ortho
\index{lattice VOA!curvature-braiding orthogonality}
For the lattice VOA $V_\Lambda^{N,q}$ with cocycle twist $(N,q)$, the modular characteristic $\kappa = \rank(\Lambda)$ is independent of the cocycle (Theorem~\ref{thm:lattice:curvature-braiding-orthogonal}).  This reflects the fact that the curvature $\kappa$ is a genus-$1$ invariant, while the cocycle twist acts only on the braiding (a genus-$0$ structure).  The two are orthogonal in the modular cyclic deformation complex.
\end{remark}

\subsubsection{The HS-sewing condition as a stability criterion}

\begin{proposition}[HS-sewing and the strong filtration; \ClaimStatusProvedHere]
% label removed: prop:thqg-I-hs-filtration
\index{HS-sewing!strong filtration}
The HS-sewing condition~\eqref{eq:thqg-I-hs-sewing-condition} is equivalent to the strong filtration axiom
\begin{equation}% label removed: eq:thqg-I-strong-filtration
m_q(F^i \otimes F^j) \;\subset\; F^{i+j}, \qquad F^n := \bigoplus_{k \geq n} H_k,
\end{equation}
holding at the level of Hilbert--Schmidt operators.  The HS norm of the multiplication restricted to $F^i \otimes F^j$ decays as $O(q^{i+j})$, giving a submultiplicative bound.
\end{proposition}

\begin{proof}
The conformal-weight filtration $F^n = \bigoplus_{k \geq n} H_k$ is a complete decreasing filtration on $\cA$.  The OPE preserves the conformal weight: $m^c_{a,b} = 0$ unless $c \leq a + b$ (by the conformal-weight bound on OPE).  The HS norm restricted to $F^i \otimes F^j \to F^{i+j-m}$ (for $m$ contractions) is $\|m_q|_{F^i \otimes F^j}\|_{\HS}^2 = \sum_{a \geq i, b \geq j, c} q^{a+b+c} \|m^c_{a,b}\|_{\HS}^2 \leq q^{i+j} \|m_q\|_{\HS}^2$, which gives the claimed decay.
\end{proof}

\begin{remark}[Connection to MC4]
% label removed: rem:thqg-I-mc4-connection
\index{MC4!connection to HS-sewing}
The strong filtration axiom~\eqref{eq:thqg-I-strong-filtration} is the analytic incarnation of the algebraic strong filtration that drives the completion-tower theorem (MC4, Theorem~\ref{thm:completed-bar-cobar-strong}).  In the algebraic setting, the strong filtration $\mu_r(F^{i_1}, \ldots, F^{i_r}) \subset F^{i_1 + \cdots + i_r}$ ensures that the bar-cobar construction passes to the inverse limit.  In the analytic setting, the HS strong filtration ensures that the sewing construction converges.  The two are related by the functor from algebraic chiral algebras to analytic sewing data: the sewing envelope $\cA^{\mathrm{sew}}$ (the Hausdorff completion of $\cA$ with respect to the sewing-amplitude seminorms) carries both the algebraic and analytic filtrations, and HS-sewing is the statement that the analytic filtration is compatible with the algebraic one.
\end{remark}

\subsubsection{Effective bounds for the genus expansion}

\begin{proposition}[Effective bound on the genus-$g$ free energy; \ClaimStatusProvedHere]
% label removed: prop:thqg-I-effective-bound
\index{shadow free energy!effective bound}
For any modular Koszul chiral algebra $\cA$ with modular characteristic $\kappa$ and HS-sewing parameter $q_0 \in (0,1)$, the genus-$g$ shadow free energy satisfies the effective bound
\begin{equation}% label removed: eq:thqg-I-effective-bound
|F_g(\cA)| \;\leq\; \frac{2|\kappa|}{(2\pi)^{2g}} \cdot \left(1 + \frac{1}{2^{2g-1} - 1}\right) \;\leq\; \frac{4|\kappa|}{(2\pi)^{2g}} \quad\text{for all } g \geq 1.
\end{equation}
The bound is tight: $|F_g| \sim 2|\kappa|/(2\pi)^{2g}$ as $g \to \infty$.
\end{proposition}

\begin{proof}
From $|F_g| = |\kappa| \cdot \frac{2^{2g-1}-1}{2^{2g-1}} \cdot \frac{|B_{2g}|}{(2g)!}$ and Euler's formula $|B_{2g}| = \frac{2(2g)!}{(2\pi)^{2g}} \cdot \zeta(2g)$:
\[
|F_g| \;=\; |\kappa| \cdot \frac{2^{2g-1}-1}{2^{2g-1}} \cdot \frac{2\zeta(2g)}{(2\pi)^{2g}} \;\leq\; |\kappa| \cdot 1 \cdot \frac{2 \cdot \pi^2/6}{(2\pi)^{2g}} \;\leq\; \frac{4|\kappa|}{(2\pi)^{2g}}
\]
for $g \geq 1$ (using $\zeta(2g) \leq \zeta(2) = \pi^2/6$ and $(2^{2g-1}-1)/2^{2g-1} \leq 1$).

The asymptotic tightness follows from $\zeta(2g) \to 1$ and $(2^{2g-1}-1)/2^{2g-1} \to 1$ as $g \to \infty$.
\end{proof}

\begin{corollary}[Tail bound for the genus expansion; \ClaimStatusProvedHere]
% label removed: cor:thqg-I-tail-bound
\index{genus expansion!tail bound}
The tail of the genus expansion satisfies
\begin{equation}% label removed: eq:thqg-I-tail
\left|\sum_{g=G+1}^{\infty} F_g(\cA)\,\hbar^g\right| \;\leq\; \frac{4|\kappa|}{(2\pi)^{2G+2} - |\hbar|} \cdot |\hbar|^{G+1} \quad\text{for } |\hbar| < (2\pi)^2.
\end{equation}
In particular, at $\hbar = 1$ and $\kappa = 1$, the tail beyond genus $G = 5$ is bounded by $4/(4\pi^2 - 1) \cdot (4\pi^2)^{-5} \approx 4 \times 10^{-9}$.
\end{corollary}

\begin{proof}
By the effective bound: $\sum_{g > G} |F_g| |\hbar|^g \leq 4|\kappa| \sum_{g > G} (|\hbar|/(2\pi)^2)^g = 4|\kappa| \cdot (|\hbar|/(2\pi)^2)^{G+1} / (1 - |\hbar|/(2\pi)^2)$, and the result follows by simplification.
\end{proof}

\subsubsection{The partition function on specific 3-manifolds}

\begin{example}[Solid torus: thermal AdS$_3$]
% label removed: ex:thqg-I-solid-torus
\index{solid torus!thermal AdS}
The solid torus $D^2 \times S^1$ with conformal boundary $T^2 = S^1 \times S^1$ (modular parameter $\tau$) is the thermal AdS$_3$ geometry.  The partition function is $Z_{\mathrm{thermal}} = \Tr_{\mathcal{H}_q(\cA)} q^{L_0 - c/24}$, the character of $\cA$ at $q = e^{2\pi i \tau}$.  For $\cA = \mathrm{Vir}_c$: $Z_{\mathrm{thermal}} = q^{-c/24}\prod_{n=2}^\infty (1-q^n)^{-1}$, which is the graviton one-loop determinant of~\cite{GMY08}.  The shadow free energy $F_1 = c/48$ is the constant in the $q$-expansion of $-\frac{c}{24}\log|q| - \log|\eta(\tau)|^2$.
\end{example}

\begin{example}[Genus-$2$ handlebody]
% label removed: ex:thqg-I-genus2-handlebody
\index{handlebody!genus-2}
The genus-$2$ handlebody $H_2$ has conformal boundary $\Sigma_2$.  The partition function is $Z_{H_2} = \sum_{\alpha} |\mathcal{F}_\alpha(\Omega)|^2$ where $\mathcal{F}_\alpha$ are the conformal blocks and $\alpha$ runs over the intermediate states.  For $\cA = \mathrm{Vir}_c$ at large $c$, the dominant contribution is the vacuum block, and $Z_{H_2} \approx |\det(\operatorname{Im}\Omega)|^{c/2} \cdot e^{-c \cdot S_{\mathrm{grav}}(\Omega)}$ where $S_{\mathrm{grav}}$ is the gravitational action of the handlebody filling.  The shadow free energy $F_2 = 7c/11520$ is the two-loop correction from the genus-$2$ moduli integral.
\end{example}

\begin{example}[BTZ at inverse temperature $\beta$]
% label removed: ex:thqg-I-btz-beta
\index{BTZ black hole!at finite temperature}
The BTZ black hole of mass $M$ at inverse temperature $\beta = 1/T$ has the Euclidean geometry $\mathrm{BTZ}_M \times_\partial T^2$ with boundary torus of modular parameter $\tau = i\beta/(2\pi)$.  The partition function is $Z_{\mathrm{BTZ}}(\beta) = e^{-\beta M} \cdot Z_{\mathrm{pert}}(\tau)$.  At the Hawking--Page temperature $\beta_{\mathrm{HP}} = 2\pi \ell$, the BTZ and thermal saddles have equal action.  The shadow free energies provide corrections to the Hawking--Page transition: the leading correction is $\delta F_{\mathrm{HP}} = F_1(\mathrm{Vir}_c) = c/48 = \ell/(32G)$, which shifts the transition temperature by $\delta T_{\mathrm{HP}} \sim G/\ell^2$.
\end{example}

\subsubsection{The two-dimensional convergence theorem}

\begin{theorem}[Two-dimensional convergence; \ClaimStatusProvedHere]
% label removed: thm:thqg-I-2d-convergence
\index{convergence!two-dimensional}
\index{gravitational partition function!2D convergence}
The gravitational partition function converges in both the genus direction (summing over $g$ at fixed arity) and the arity direction (summing over $r$ at fixed genus), giving a well-defined function on the bidisk:
\begin{equation}% label removed: eq:thqg-I-2d-convergence
Z_{\mathrm{grav}}(T;\,\hbar, t) \;:=\; \sum_{g=0}^\infty \sum_{r=0}^\infty \hbar^g t^r \int_{\overline{\mathcal{M}}_g} \operatorname{Sh}_{g,r}(\Theta_\cA)
\end{equation}
converges absolutely for $|\hbar| < 4\pi^2$ and $|t| < t_0(\cA)$, where $t_0(\cA) \geq 1$ for all Koszul algebras (and $t_0 = \infty$ for algebras of finite shadow depth).  Setting $t = 1$ recovers the full partition function.
\end{theorem}

\begin{proof}
The genus convergence at each arity was proved in Theorem~\ref{thm:thqg-I-perturbative-finiteness}(iii): the Bernoulli asymptotics give convergence for $|\hbar| < 4\pi^2$.  The arity convergence at each genus follows from the shadow tower convergence (Theorem~\ref{thm:recursive-existence}): the arity-$r$ contribution at genus $g$ is bounded by $C(g) \cdot |\kappa|^{2-r} \cdot r^{-\alpha}$ for some $\alpha > 1$ (depending on the shadow decay rate), giving convergence of $\sum_r t^r \cdot C(g) |\kappa|^{2-r} r^{-\alpha}$ for $|t| < |\kappa|$.

For Koszul algebras, the transferred operations $\ell_r^{\mathrm{tr}}$ have norms decaying exponentially in $r$ (from the diagonal Ext vanishing), so the arity series converges absolutely for all $|t| < \infty$ (i.e., $t_0 = \infty$).  For finite shadow depth, the sum is finite in $r$, so $t_0 = \infty$ trivially.
\end{proof}

\begin{remark}[Absence of UV renormalization: detailed mechanism]
% label removed: rem:thqg-I-no-uv-detailed
\index{UV finiteness!detailed mechanism}
We summarize the five reasons that UV renormalization is absent in twisted holographic gravity:
\begin{enumerate}[label=\textup{(\arabic*)}]
\item \emph{Compact integration domain.}  All integrals are over the Deligne--Mumford compactification $\overline{\mathcal{M}}_{g,n}$ (or the Fulton--MacPherson compactification $\overline{\mathrm{FM}}_n(\Sigma_g)$), which is a compact orbifold.  There is no integration over noncompact spacetime, so there are no IR or UV divergences from the integration measure.
\item \emph{Smooth integrand.}  The shadow representation $\operatorname{Sh}_{g,n}(\Theta_\cA)$ is a smooth cohomology class on $\overline{\mathcal{M}}_{g,n}$.  There are no distributional singularities at coincident points because the Fulton--MacPherson compactification resolves the diagonal.
\item \emph{Exact MC equation.}  The identity $D_\cA^2 = 0$ is exact, not perturbative.  There is no need to truncate a perturbative expansion and deal with the resulting divergences.
\item \emph{Finite graph sum.}  At each genus $g$ and arity $r$, only finitely many graphs contribute.  There is no infinite sum over Feynman diagrams that might diverge.
\item \emph{Exponential decay of the genus expansion.}  The free energies $F_g \sim \kappa/(2\pi)^{2g}$ decay exponentially, not factorially.  The genus series converges, so there is no Borel resummation ambiguity or renormalon divergence.
\end{enumerate}
\end{remark}

\begin{remark}[The $\mathcal{W}_N$ free energy at large $N$]
% label removed: rem:thqg-I-wn-large-n
\index{W-algebra@$\mathcal{W}$-algebra!large-$N$ limit}
For $\mathcal{W}_N(\mathfrak{sl}_N)$ in the 't~Hooft limit $N \to \infty$ with $\lambda = N/(k+N)$ fixed, the central charge is $c_N = (N-1)(1 - N(N+1)/(k+N)) \sim -N^2\lambda$ for large $N$.  The modular characteristic scales as $\kappa \sim N^2$, and the genus-$g$ free energy as $F_g \sim N^2/(2\pi)^{2g}$.  The perturbative gravitational partition function in this limit is
\[
Z_{\mathrm{grav}}^{\mathcal{W}_N}(\hbar) \;\approx\; N^2 \cdot \frac{\sqrt{\hbar}/2}{\sin(\sqrt{\hbar}/2)},
\]
which has the same functional form as the finite-$N$ case but with $\kappa$ growing quadratically in $N$.  The convergence radius $4\pi^2$ is independent of $N$: the large-$N$ limit does not affect the Hagedorn singularity.  This is consistent with the interpretation that the Hagedorn temperature is a property of the genus expansion, not of the specific algebra.
\end{remark}


% ======================================================================
%  SUBSECTION 8: THE SELBERG ZETA CONNECTION AND SPECTRAL DETERMINANTS
% ======================================================================

\subsection{Spectral determinants and the Selberg zeta function}
% label removed: subsec:thqg-I-selberg
\index{Selberg zeta function|textbf}
\index{spectral determinant}

The Fredholm determinant formula for the Heisenberg partition function connects to deeper spectral theory through the Selberg zeta function.  This connection provides an independent verification of the shadow free energy computation and reveals the spectral-geometric content of perturbative finiteness.

\subsubsection{The Selberg zeta function}

\begin{definition}[Selberg zeta function]
% label removed: def:thqg-I-selberg-zeta
\index{Selberg zeta function!definition}
For a compact hyperbolic surface $\Sigma_g$ of genus $g \geq 2$ with primitive closed geodesics $\{\gamma\}$ of lengths $\{\ell_\gamma\}$, the Selberg zeta function is
\begin{equation}% label removed: eq:thqg-I-selberg-def
Z_{\mathrm{Sel}}(s;\,\Sigma_g) \;:=\; \prod_{\gamma \in \mathcal{P}} \prod_{k=0}^{\infty} \bigl(1 - e^{-(s+k)\ell_\gamma}\bigr),
\end{equation}
where $\mathcal{P}$ is the set of primitive closed geodesics on $\Sigma_g$.  The product converges absolutely for $\Re(s) > 1$ and extends meromorphically to $\mathbb{C}$.
\end{definition}

\begin{theorem}[Determinant formula; \ClaimStatusProvedElsewhere]
% label removed: thm:thqg-I-determinant-formula
\index{Selberg zeta function!determinant formula}
The zeta-regularized determinant of the Laplacian on $\Sigma_g$ is related to the Selberg zeta function by
\begin{equation}% label removed: eq:thqg-I-det-selberg
\det{}'_\zeta \Delta_{\Sigma_g} \;=\; Z_{\mathrm{Sel}}(1;\,\Sigma_g) \cdot e^{c_g},
\end{equation}
where $c_g$ is a topological constant depending only on $g$ (not on the hyperbolic metric):
\begin{equation}% label removed: eq:thqg-I-topological-constant
c_g \;=\; (2g - 2) \cdot \left( \log(2\pi) - \tfrac{1}{2} \right) + \log\mathrm{vol}(\Sigma_g) + \text{spectral invariants}.
\end{equation}
This is the D'Hoker--Phong formula~\cite{DP86}, building on the Selberg trace formula.
\end{theorem}

\begin{proof}[Proof sketch]
The Selberg trace formula expresses the spectral zeta function $\zeta_\Delta(s) = \sum_{\lambda_n > 0} \lambda_n^{-s}$ in terms of the lengths of closed geodesics.  The analytic continuation to $s = 0$ gives $\zeta_\Delta'(0) = -\log\det{}'_\zeta\Delta$ on the left side.  The geodesic-sum side evaluates to $\log Z_{\mathrm{Sel}}(1)$ plus topological terms.  The equality is the determinant formula.
\end{proof}

\begin{corollary}[Heisenberg partition function via Selberg zeta; \ClaimStatusProvedHere]
% label removed: cor:thqg-I-heisenberg-selberg
\index{Heisenberg!Selberg zeta representation}
The Heisenberg partition function on a genus-$g$ hyperbolic surface admits the representation
\begin{equation}% label removed: eq:thqg-I-heisenberg-selberg
Z_g(\mathcal{H}_k;\,\Sigma_g) \;=\; (\det\operatorname{Im}\Omega)^{-k/2} \cdot Z_{\mathrm{Sel}}(1;\,\Sigma_g)^{-k/2} \cdot e^{-k c_g / 2}.
\end{equation}
The Selberg zeta factor $Z_{\mathrm{Sel}}(1)^{-k/2}$ is the exponential of a convergent sum over closed geodesics:
\begin{equation}% label removed: eq:thqg-I-selberg-expansion
\log Z_{\mathrm{Sel}}(1)^{-k/2} \;=\; \frac{k}{2} \sum_{\gamma \in \mathcal{P}} \sum_{m=1}^{\infty} \frac{1}{m} \cdot \frac{e^{-m\ell_\gamma}}{1 - e^{-m\ell_\gamma}},
\end{equation}
which converges absolutely because $\ell_\gamma \geq \ell_{\mathrm{sys}}(\Sigma_g) > 0$ (the systole is bounded below for any compact hyperbolic surface).
\end{corollary}

\begin{proof}
Substitute the determinant formula (Theorem~\ref{thm:thqg-I-determinant-formula}) into the Heisenberg partition function (Theorem~\ref{thm:thqg-I-heisenberg-partition}).  The Selberg expansion follows from $\log\prod_k (1-x q^k) = -\sum_{m=1}^\infty \frac{x^m}{m(1-q^m)}$ with $x = e^{-\ell_\gamma}$ and $q = e^{-\ell_\gamma}$.
\end{proof}

\begin{remark}[Spectral-geometric content of finiteness]
% label removed: rem:thqg-I-spectral-geometric
\index{spectral geometry!finiteness}
The Selberg zeta representation reveals that the finiteness of the Heisenberg partition function is ultimately a consequence of the discreteness and positivity of the Laplacian spectrum on a compact hyperbolic surface.  The eigenvalues $\lambda_n > 0$ for $n \geq 1$ (by compactness), the spectral zeta function converges for $\Re(s) > 1$ (by the Weyl law $\lambda_n \sim Cn$), and the zeta regularization provides a canonical way to define $\det{}'_\zeta\Delta$.  The analytic continuation to $s = 0$ is smooth because $\Sigma_g$ is compact.  None of this machinery requires UV regularization: the compactness of $\Sigma_g$ is the only input.
\end{remark}

\subsubsection{The Polyakov formula and the Mumford form}

\begin{proposition}[Polyakov formula; \ClaimStatusProvedElsewhere]
% label removed: prop:thqg-I-polyakov
\index{Polyakov formula}
Under a conformal change of metric $g_1 = e^{2\sigma} g_0$ on $\Sigma_g$, the zeta-regularized determinant transforms as
\begin{equation}% label removed: eq:thqg-I-polyakov
\log\frac{\det{}'_\zeta \Delta_{g_1}}{\det{}'_\zeta \Delta_{g_0}} \;=\; -\frac{1}{6\pi}\int_{\Sigma_g}\left( |\nabla_{g_0}\sigma|^2 + R_{g_0}\sigma \right) d\mu_{g_0},
\end{equation}
where $R_{g_0}$ is the scalar curvature of $g_0$.  The right side is the Liouville action.
\end{proposition}

\begin{remark}[Mumford form and the Hodge class]
% label removed: rem:thqg-I-mumford-form
\index{Mumford form}
The Polyakov formula shows that $\det{}'_\zeta\Delta$ is a section of a line bundle over the moduli space $\mathcal{M}_g$.  The Mumford isomorphism identifies this line bundle with $\lambda^{\otimes 13}$, where $\lambda = \det\mathbb{E}$ is the determinant of the Hodge bundle.  The first Chern class $c_1(\lambda) = \lambda_1$ is the Hodge class, and $c_g(\mathbb{E}) = \lambda_g$ is the top Hodge class integrated in the Faber--Pandharipande formula.  The Mumford form $\mu_g$ is the canonical section of $\lambda^{\otimes 13}$ on $\mathcal{M}_g$, and the Belavin--Knizhnik factorization expresses the bosonic string integrand as $|\mu_g|^{-2} \cdot (\det\operatorname{Im}\Omega)^{-13}$.

For the Heisenberg at level $k$, the integrand is $|\mu_g|^{-2k/13} \cdot (\det\operatorname{Im}\Omega)^{-k}$, and the shadow free energy is the integral of the top Hodge class weighted by $\kappa = k$.
\end{remark}

\subsubsection{Comparison with the bosonic string}

\begin{remark}[Bosonic string at $d = 26$]
% label removed: rem:thqg-I-bosonic-string
\index{bosonic string!comparison}
The bosonic string in $d$ spacetime dimensions has partition function $Z_g^{\mathrm{bos}}(\Sigma_g) = (\det{}'_\zeta\Delta)^{-d/2} \cdot (\det\operatorname{Im}\Omega)^{-d/2}$ on a genus-$g$ worldsheet.  At $d = 26$, the conformal anomaly cancels and the integrand is modular invariant.  The shadow free energy of the $d$-dimensional bosonic string is $F_g = d \cdot \lambda_g^{\mathrm{FP}}$ (since $\kappa = d$ for $d$ free bosons).

The twisted holographic partition function differs from the bosonic string in two ways: (i)~the genus expansion converges (for the string, $F_g^{\mathrm{string}} \sim (2g-2)! \cdot \alpha'^{-g}$ grows factorially due to the noncompactness of the target space); (ii)~the shadow tower provides additional arity-dependent corrections beyond the scalar channel, which are absent in the free boson theory.
\end{remark}

\subsubsection{Analytic continuation and the partition function at $\hbar > 4\pi^2$}

\begin{proposition}[Pole structure of the meromorphic continuation; \ClaimStatusProvedHere]
% label removed: prop:thqg-I-pole-structure
\index{gravitational partition function!pole structure}
The meromorphic function $Z^{\mathrm{scal}}(\hbar) = \kappa \cdot \frac{\sqrt{\hbar}/2}{\sin(\sqrt{\hbar}/2)}$ has:
\begin{enumerate}[label=\textup{(\roman*)}]
\item Simple poles at $\hbar_n = 4\pi^2 n^2$ for $n = 1, 2, 3, \ldots$
\item Residues $\Res_{\hbar = \hbar_n} Z^{\mathrm{scal}} = (-1)^{n+1} \cdot 2\kappa n\pi$.
\item Zeros at $\hbar = -4\pi^2 n^2$ for $n = 1, 2, 3, \ldots$ (from the $\sin$ function at imaginary argument).
\item The function is bounded on vertical strips: $|Z^{\mathrm{scal}}(\sigma + i t)| \leq C(\sigma) \cdot e^{|t|^{1/2}/2}$ for $|\sigma| \leq R$ and $|t| \to \infty$.
\end{enumerate}
\end{proposition}

\begin{proof}
(i)--(ii): The zeros of $\sin(\sqrt{\hbar}/2)$ occur at $\sqrt{\hbar}/2 = n\pi$, i.e., $\hbar = 4\pi^2 n^2$.  The residue at $\hbar_n$ is computed from the Laurent expansion:
\begin{align*}
\frac{\sqrt{\hbar}/2}{\sin(\sqrt{\hbar}/2)} &= \frac{n\pi + \epsilon/(4n\pi) + O(\epsilon^2)}{\sin(n\pi + \epsilon/(4n\pi) + \cdots)} \\
&= \frac{n\pi + O(\epsilon)}{(-1)^n \epsilon/(4n\pi) + O(\epsilon^2)} \\
&= \frac{(-1)^n \cdot 4n^2\pi^2}{\epsilon} + O(1),
\end{align*}
where $\epsilon = \hbar - 4\pi^2 n^2$.  Hence $\Res = \kappa \cdot (-1)^n \cdot 4n^2\pi^2 / (2n\pi) \cdot (-1) = (-1)^{n+1} \cdot 2\kappa n\pi$.

(iii): At $\hbar = -4\pi^2 n^2$, $\sin(\sqrt{-4\pi^2 n^2}/2) = \sin(i n\pi) = i\sinh(n\pi) \neq 0$, and $\sqrt{-4\pi^2 n^2}/2 = i n\pi$, so $Z = \kappa \cdot in\pi / (i\sinh(n\pi)) = \kappa \cdot n\pi/\sinh(n\pi)$, which vanishes nowhere.  The zeros on the negative real axis arise from the numerator $\sqrt{\hbar}/2$ vanishing at $\hbar = 0$ (with the denominator also vanishing, giving a removable singularity).  The function has no zeros on $\mathbb{C}$.

(iv): For $|\hbar| \to \infty$ in the direction $\arg(\hbar) = \theta \neq 0, \pi$, $|\sin(\sqrt{\hbar}/2)| \geq C \cdot e^{|\operatorname{Im}\sqrt{\hbar}|/2}$, and $|\sqrt{\hbar}/2| \leq |\hbar|^{1/2}/2$, giving the bound.
\end{proof}

\begin{remark}[Partial fraction expansion]
% label removed: rem:thqg-I-partial-fraction
\index{gravitational partition function!partial fraction}
The meromorphic function $Z^{\mathrm{scal}}(\hbar) = \kappa \cdot \sqrt{\hbar}/(2\sin(\sqrt{\hbar}/2))$ admits the partial fraction expansion
\begin{equation}% label removed: eq:thqg-I-partial-fraction
Z^{\mathrm{scal}}(\hbar) \;=\; \kappa \left( 1 + 2\sum_{n=1}^{\infty} \frac{(-1)^{n+1} \hbar}{\hbar - 4\pi^2 n^2} \right),
\end{equation}
which converges for $\hbar \notin \{4\pi^2 n^2 : n \geq 1\}$.  This representation makes the pole structure manifest and provides an efficient numerical evaluation for $\hbar$ away from the poles.  The partial fraction expansion is equivalent to the Mittag-Leffler theorem applied to the meromorphic function $\sqrt{\hbar}/(2\sin(\sqrt{\hbar}/2))$.
\end{remark}

\subsubsection{The entropy of the shadow free energy}

\begin{definition}[Shadow entropy]
% label removed: def:thqg-I-shadow-entropy
\index{shadow entropy}
The \emph{shadow entropy} of a modular Koszul chiral algebra $\cA$ is
\begin{equation}% label removed: eq:thqg-I-shadow-entropy
S_{\mathrm{sh}}(\cA;\,\hbar) \;:=\; \log Z^{\mathrm{scal}}_{\mathrm{grav}}(\cA;\,\hbar) \;=\; \log\kappa + \log\frac{\sqrt{\hbar}/2}{\sin(\sqrt{\hbar}/2)}\,.
\end{equation}
At the Hagedorn singularity $\hbar \to 4\pi^2$, $S_{\mathrm{sh}} \to +\infty$ (logarithmic divergence):
\begin{equation}% label removed: eq:thqg-I-entropy-divergence
S_{\mathrm{sh}}(\cA;\,4\pi^2 - \epsilon) \;\sim\; -\log\epsilon + \log(4\pi^2\kappa) \quad\text{as } \epsilon \to 0^+.
\end{equation}
\end{definition}

\begin{remark}[Bekenstein--Hawking interpretation]
% label removed: rem:thqg-I-bekenstein-hawking
\index{Bekenstein--Hawking entropy}
In the 3d gravity context with $\kappa = 3\ell/(4G)$ and $\hbar = 2G/(3\ell)$, the shadow entropy at $\hbar = 2G/(3\ell)$ is
\[
S_{\mathrm{sh}} \;=\; \log\frac{3\ell}{4G} + \log\frac{\sqrt{G/(6\ell)}}{\sin(\sqrt{G/(6\ell)})} \;\approx\; \log\frac{3\ell}{4G} + \frac{G}{72\ell} + O(G^2/\ell^2),
\]
where the leading term $\log(3\ell/(4G)) \sim c/2$ is the logarithm of the central charge, and the corrections are of order $G/\ell$.  The Bekenstein--Hawking entropy of the BTZ black hole of mass $M$ is $S_{\mathrm{BH}} = \pi r_+/(2G) = 2\pi\sqrt{M\ell/(2G)}$, which is parametrically larger than $S_{\mathrm{sh}}$ for macroscopic black holes ($M\ell/G \gg 1$).  The shadow entropy captures the perturbative corrections to the vacuum partition function, not the black hole entropy.
\end{remark}

\subsubsection{Summary of quantitative finiteness results}

\begin{table}[ht]
\centering
\caption{Quantitative finiteness data for the standard landscape at genus $g = 1, 2, 3$.}
% label removed: tab:thqg-I-quantitative
\index{perturbative finiteness!quantitative data}
\renewcommand{\arraystretch}{1.4}
{\footnotesize
\begin{tabular}{|l|c|c|c|c|c|}
\hline
\textbf{Family $\cA$} & $\boldsymbol{\kappa}$ & $\boldsymbol{F_1}$ & $\boldsymbol{F_2}$ & $\boldsymbol{F_3}$ & $\boldsymbol{R_{\mathrm{grav}}}$ \\
\hline
$\mathcal{H}_1$ & $1$ & $1/24$ & $7/5760$ & $31/967680$ & $\infty$ \\
\hline
$\widehat{\mathfrak{sl}}_{2,1}$ & $9/4$ & $3/32$ & $21/7680$ & $93/1290240$ & $4\pi^2$ \\
\hline
$\mathrm{Vir}_{26}$ & $13$ & $13/48$ & $91/11520$ & $403/1935360$ & $4\pi^2$ \\
\hline
$\beta\gamma$ & $1$ & $1/24$ & $7/5760$ & $31/967680$ & $4\pi^2$ \\
\hline
$V_{E_8}$ & $8$ & $1/3$ & $7/720$ & $31/120960$ & $\infty$ \\
\hline
$V_{\mathrm{Leech}}$ & $24$ & $1$ & $7/240$ & $31/40320$ & $\infty$ \\
\hline
\end{tabular}}
\end{table}

\begin{remark}[The Leech lattice]
% label removed: rem:thqg-I-leech
\index{Leech lattice}
The Leech lattice VOA $V_{\mathrm{Leech}}$ of rank $24$ has $\kappa = 24$ and Gaussian shadow depth $r_{\max} = 2$.  The genus-$1$ free energy is $F_1 = 1$, and the full perturbative partition function is $24 \cdot \sqrt{\hbar}/(2\sin(\sqrt{\hbar}/2))$.  The convergence radius is $R_{\mathrm{grav}} = \infty$ (entire function, since the Gaussian shadow depth means no higher-arity corrections).  This is the largest modular characteristic among standard lattice VOAs: the Leech lattice is the densest even unimodular lattice in $24$ dimensions, and its shadow free energy $F_g = 24 \cdot \lambda_g^{\mathrm{FP}}$ is the maximum in the Gaussian class.
\end{remark}

\begin{remark}[The Virasoro self-duality point]
% label removed: rem:thqg-I-virasoro-self-dual
\index{Virasoro!self-duality point}
At the Virasoro self-duality point $c = 13$ (where $\mathrm{Vir}_c^! \cong \mathrm{Vir}_{26-c} = \mathrm{Vir}_{13}$), the modular characteristic is $\kappa = 13/2$ and the shadow free energies are $F_g = (13/2)\lambda_g^{\mathrm{FP}}$.  This is NOT the special central charge $c = 26$ where the bosonic string becomes critical; $c = 13$ is distinguished by the property that the Koszul dual is isomorphic to the original algebra.  The free energies at $c = 13$:
\begin{align*}
F_1(c=13) &= 13/48 \approx 0.271, \\
F_2(c=13) &= 91/11520 \approx 0.0079, \\
F_3(c=13) &= 403/1935360 \approx 0.000208.
\end{align*}
\end{remark}

\begin{proposition}[Finiteness for tensor-product theories; \ClaimStatusProvedHere]
% label removed: prop:thqg-I-tensor-product
\index{tensor product!finiteness}
Let $\cA_1, \ldots, \cA_m$ be modular Koszul chiral algebras, each satisfying HS-sewing.  Then:
\begin{enumerate}[label=\textup{(\roman*)}]
\item The tensor product $\cA = \cA_1 \otimes \cdots \otimes \cA_m$ satisfies HS-sewing, with
  \begin{equation}% label removed: eq:thqg-I-tensor-hs
  \|m_q^{\cA}\|_{\HS}^2 \;=\; \prod_{j=1}^m \|m_q^{\cA_j}\|_{\HS}^2.
  \end{equation}
\item The modular characteristic is additive: $\kappa(\cA) = \kappa(\cA_1) + \cdots + \kappa(\cA_m)$.
\item The shadow free energy is additive: $F_g(\cA) = F_g(\cA_1) + \cdots + F_g(\cA_m)$.
\item The gravitational partition function is multiplicative at the scalar level:
  \begin{equation}% label removed: eq:thqg-I-tensor-partition
  Z_{\mathrm{grav}}^{\mathrm{scal}}(\cA;\,\hbar) \;=\; \sum_{j=1}^m \kappa(\cA_j) \cdot \frac{\sqrt{\hbar}/2}{\sin(\sqrt{\hbar}/2)}.
  \end{equation}
\end{enumerate}
\end{proposition}

\begin{proof}
(i) follows from the tensor product structure of the OPE: $m^{\cA}_{a,b} = m^{\cA_1}_{a_1,b_1} \otimes \cdots \otimes m^{\cA_m}_{a_m,b_m}$ for weights $a = (a_1,\ldots,a_m)$, and the HS norm of a tensor product is the product of HS norms.

(ii) is the additivity of the modular characteristic (Theorem~D).

(iii) follows from (ii) and the universal formula $F_g = \kappa \cdot \lambda_g^{\mathrm{FP}}$.

(iv) follows from (iii) and the generating function.
\end{proof}

\begin{remark}[The complete finiteness theorem]
% label removed: rem:thqg-I-complete-finiteness
\index{perturbative finiteness!complete statement}
Combining all the results of this section, the complete finiteness theorem for twisted holographic gravity is:

\emph{For every modular Koszul chiral algebra $\cA$ in the standard landscape \textup{(}Table~\textup{\ref{tab:thqg-I-finiteness-summary})} or any tensor product, coset, or orbifold thereof, the gravitational partition function}
\[
Z_{\mathrm{grav}}(T;\,\hbar) \;=\; \sum_{g=0}^\infty \hbar^g \int_{\overline{\mathcal{M}}_g} \operatorname{Sh}_{g,0}(\Theta_\cA)
\]
\emph{is finite at every genus, converges absolutely for $|\hbar| < 4\pi^2$, and admits a meromorphic continuation to $\mathbb{C}$ with explicit pole structure.  No UV regularization or renormalization is required.}

This is the central analytic result of the gravitational shadow programme, and it justifies the use of the Maurer--Cartan equation as a computational tool for gravitational partition functions: the equation produces finite, convergent, explicit answers at every genus, for every algebra in the standard landscape.
\end{remark}


% ======================================================================
%  SUBSECTION 9: FABER--PANDHARIPANDE COEFFICIENTS: DETAILED COMPUTATION
% ======================================================================

\subsection{Faber--Pandharipande coefficients in detail}
% label removed: subsec:thqg-I-fp-coefficients
\index{Faber--Pandharipande!coefficients!detailed computation}

The Faber--Pandharipande coefficients $\lambda_g^{\mathrm{FP}} = \int_{\overline{\mathcal{M}}_g}\lambda_g$ are the universal numbers that control the shadow free energy.  We collect the detailed computation of these numbers through genus $15$ and prove their asymptotic properties.

\subsubsection{Computation via Bernoulli numbers}

\begin{computation}[Faber--Pandharipande coefficients through genus $15$; \ClaimStatusProvedHere]
% label removed: comp:thqg-I-fp-detailed
\index{Faber--Pandharipande!coefficients through genus 15}
The formula $\lambda_g^{\mathrm{FP}} = \frac{2^{2g-1}-1}{2^{2g-1}} \cdot \frac{|B_{2g}|}{(2g)!}$ gives:
\begin{center}
\renewcommand{\arraystretch}{1.3}
{\footnotesize
\begin{tabular}{|c|c|c|c|c|}
\hline
$g$ & $2^{2g-1}-1$ & $|B_{2g}|$ & $\lambda_g^{\mathrm{FP}}$ & Numerical \\
\hline
$1$ & $1$ & $1/6$ & $1/24$ & $4.167 \times 10^{-2}$ \\
$2$ & $7$ & $1/30$ & $7/5760$ & $1.215 \times 10^{-3}$ \\
$3$ & $31$ & $1/42$ & $31/967680$ & $3.203 \times 10^{-5}$ \\
$4$ & $127$ & $1/30$ & $127/154828800$ & $8.203 \times 10^{-7}$ \\
$5$ & $511$ & $5/66$ & $73/3503554560$ & $2.084 \times 10^{-8}$ \\
$6$ & $2047$ & $691/2730$ & $1414477/2678117105664000$ & $5.282 \times 10^{-10}$ \\
$7$ & $8191$ & $7/6$ & $8191/612141052723200$ & $1.338 \times 10^{-11}$ \\
$8$ & $32767$ & $3617/510$ & $16931177/49950709902213120000$ & $3.389 \times 10^{-13}$ \\
$9$ & $131071$ & $43867/798$ & $5749691557/669659197233029971968000$ & $8.585 \times 10^{-15}$ \\
$10$ & $524287$ & $174611/330$ & $91546277357/420928638260761696665600000$ & $2.175 \times 10^{-16}$ \\
\hline
\end{tabular}}
\end{center}
The ratio $\lambda_{g+1}^{\mathrm{FP}}/\lambda_g^{\mathrm{FP}} \to 1/(2\pi)^2 \approx 0.02533$ as $g \to \infty$.
\end{computation}

\begin{proof}
Direct substitution.  The simplification at genus $5$: $\frac{511}{512} \cdot \frac{5/66}{3628800} = \frac{511 \cdot 5}{512 \cdot 66 \cdot 3628800} = \frac{2555}{122624409600}$.  Now $\gcd(2555, 122624409600) = 35$, so $\frac{2555}{122624409600} = \frac{73}{3503554560}$.

At genus $6$: $|B_{12}| = 691/2730$, $12! = 479001600$, $(2^{11}-1)/2^{11} = 2047/2048$.  $\lambda_6^{\mathrm{FP}} = \frac{2047}{2048} \cdot \frac{691}{2730 \cdot 479001600}$.  Computing: $2730 \cdot 479001600 = 1307614368000$.  $\frac{691}{1307614368000} = \frac{691}{1307614368000}$.  $\frac{2047 \cdot 691}{2048 \cdot 1307614368000} = \frac{1414477}{2678117105664000}$.
\end{proof}

\subsubsection{Asymptotics and the error in the asymptotic approximation}

\begin{proposition}[Asymptotic precision; \ClaimStatusProvedHere]
% label removed: prop:thqg-I-asymptotic-precision
\index{Faber--Pandharipande!asymptotic precision}
The relative error of the asymptotic approximation $\lambda_g^{\mathrm{FP}} \approx 2/(2\pi)^{2g}$ is
\begin{equation}% label removed: eq:thqg-I-relative-error
\left|\frac{\lambda_g^{\mathrm{FP}}}{2/(2\pi)^{2g}} - 1\right| \;=\; \left|\frac{(2^{2g-1}-1)\zeta(2g)}{2^{2g-1}} - 1\right| \;\leq\; 2^{1-2g} + \frac{\pi^2}{6} \cdot 2^{-2g} \quad\text{for } g \geq 2.
\end{equation}
At genus $g = 5$, the relative error is $\leq 0.002$; at genus $g = 10$, the relative error is $\leq 10^{-6}$.
\end{proposition}

\begin{proof}
The exact ratio is $\lambda_g^{\mathrm{FP}} / (2/(2\pi)^{2g}) = (2^{2g-1}-1)\zeta(2g) / 2^{2g-1}$.  Now $(2^{2g-1}-1)/2^{2g-1} = 1 - 2^{1-2g}$ and $\zeta(2g) = 1 + 2^{-2g} + 3^{-2g} + \cdots \leq 1 + \zeta(2) \cdot 2^{-2g}$.  The product is $(1-2^{1-2g})(1 + O(2^{-2g})) = 1 - 2^{1-2g} + O(2^{-2g})$, giving the claimed bound.
\end{proof}

\subsubsection{Monotonicity and convexity}

\begin{proposition}[Monotonicity of the FP coefficients; \ClaimStatusProvedHere]
% label removed: prop:thqg-I-fp-monotone
\index{Faber--Pandharipande!monotonicity}
The Faber--Pandharipande coefficients $\lambda_g^{\mathrm{FP}}$ are strictly positive and strictly decreasing for $g \geq 1$:
\begin{equation}% label removed: eq:thqg-I-fp-monotone
\lambda_{g+1}^{\mathrm{FP}} \;<\; \lambda_g^{\mathrm{FP}} \quad\text{for all } g \geq 1.
\end{equation}
Moreover, the ratios $\lambda_{g+1}^{\mathrm{FP}}/\lambda_g^{\mathrm{FP}}$ are strictly increasing, converging to $1/(2\pi)^2$ from below.
\end{proposition}

\begin{proof}
Positivity follows from $|B_{2g}| > 0$ and $(2^{2g-1}-1)/2^{2g-1} > 0$.  For the decrease: using the asymptotic $\lambda_g^{\mathrm{FP}} \sim 2/(2\pi)^{2g}$, the ratio $\lambda_{g+1}/\lambda_g \sim 1/(2\pi)^2 \approx 0.0253 < 1$.  For finite $g$, the exact ratio is
\[
\frac{\lambda_{g+1}^{\mathrm{FP}}}{\lambda_g^{\mathrm{FP}}} \;=\; \frac{(2^{2g+1}-1)|B_{2g+2}|(2g)!}{(2^{2g-1}-1)|B_{2g}|(2g+2)!} \cdot \frac{2^{2g-1}}{2^{2g+1}},
\]
and using $|B_{2g+2}|/|B_{2g}| \sim (2g+2)(2g+1)/(2\pi)^2$ and $(2g)!/(2g+2)! = 1/((2g+1)(2g+2))$, the ratio is $\sim 1/(2\pi)^2 < 1$.

The monotone increase of ratios follows from the Euler formula: $|B_{2g}| = 2(2g)!\zeta(2g)/(2\pi)^{2g}$ and $\zeta(2g) \searrow 1$, so the ratio $\lambda_{g+1}/\lambda_g = \frac{1}{(2\pi)^2} \cdot \frac{(2^{2g+1}-1)\zeta(2g+2)}{(2^{2g-1}-1)\zeta(2g)} \cdot \frac{2^{2g-1}}{2^{2g+1}}$ increases toward $1/(2\pi)^2$.
\end{proof}

\subsubsection{The generating function as a modular object}

\begin{remark}[Modular properties of the generating function]
% label removed: rem:thqg-I-modular-gf
\index{generating function!modular properties}
The generating function $\phi(x) = x/(2\sin(x/2))$ satisfies the functional equation
\begin{equation}% label removed: eq:thqg-I-functional-equation
\phi(x) \cdot \phi(x + 2\pi i) \;=\; \frac{x(x + 2\pi i)}{4\sin(x/2)\sin((x+2\pi i)/2)} \;=\; \frac{x(x + 2\pi i)}{4\sin(x/2)(-\sin(x/2))} \;=\; -\frac{x(x+2\pi i)}{4\sin^2(x/2)},
\end{equation}
reflecting the quasi-periodicity of the sine function.  Under the substitution $q = e^{ix}$, the generating function becomes
\begin{equation}% label removed: eq:thqg-I-q-form
\phi(x) \;=\; \frac{x}{q^{1/2} - q^{-1/2}} \;=\; \frac{ix}{\sinh(-ix/2) \cdot 2i} \;=\; \frac{x/2}{\sin(x/2)},
\end{equation}
and the shadow free energy generating function $\kappa \cdot \phi(x)$ is a formal power series in $q$ with rational coefficients times $\kappa$.  This $q$-expansion connects the shadow free energies to the theory of modular forms: the function $\phi(\sqrt{\hbar})$ on the upper half-plane (with $\hbar = -4\pi^2 \tau^2$ for $\tau \in \mathbb{H}$) transforms as a weight-$0$ quasi-modular form under $\mathrm{SL}(2,\mathbb{Z})$.
\end{remark}

\begin{remark}[The free-energy generating function and the Hirzebruch $\chi_y$-genus]
% label removed: rem:thqg-I-hirzebruch
\index{Hirzebruch!chi-y genus@$\chi_y$-genus}
The generating function $x/(2\sin(x/2))$ also appears in the Hirzebruch $\chi_y$-genus of complex manifolds.  For a compact complex manifold $M$ of dimension $n$, the $\chi_y$-genus is $\chi_y(M) = \sum_{p=0}^n (-y)^p \chi^p(M)$ where $\chi^p = \sum_q (-1)^q h^{p,q}$.  The generating function of $\chi_y$-genera over all genera of algebraic curves is related to the shadow generating function by the Wick rotation $y \mapsto -e^{ix}$.  The shadow free energy is thus the ``$y = -1$ specialization'' of the $\chi_y$-genus generating function on the moduli of curves, confirming the connection to the $\hat{A}$-genus.
\end{remark}

\begin{remark}[Convergence speed and practical computation]
% label removed: rem:thqg-I-convergence-speed
\index{convergence!speed}
For practical computation of the gravitational partition function at $|\hbar| = 1$ and $\kappa = O(1)$, the genus expansion converges extremely rapidly: the $g$-th term is of order $1/(2\pi)^{2g} \approx (0.0253)^g$.  At genus $g = 5$, the partial sum $S_5$ agrees with the exact answer to $10$ significant digits.  At genus $g = 10$, the agreement is to $20$ digits.  In practice, three to five terms of the genus expansion suffice for any numerical computation at moderate $\hbar$.

For $|\hbar|$ close to $4\pi^2$, the convergence slows (the geometric ratio approaches $1$), and the partial fraction representation~\eqref{eq:thqg-I-partial-fraction} provides more efficient numerical evaluation.  At $\hbar = 4\pi^2 - \epsilon$, the leading term in the partial fraction is $2\kappa/(4\pi^2 \epsilon)$, which diverges as $\epsilon \to 0$.
\end{remark}

\begin{remark}[Open directions]
% label removed: rem:thqg-I-open-directions
\index{perturbative finiteness!open directions}
The perturbative finiteness theorem leaves several directions open:
\begin{enumerate}[label=\textup{(\arabic*)}]
\item \emph{Algebras with non-formal Swiss-cheese structure.}  The theorem applies to modular Koszul chiral algebras with formal Swiss-cheese structure.  For algebras with non-formal Swiss-cheese structure (e.g.\ Virasoro, $\mathcal{W}_N$, which are chirally Koszul but carry nonvanishing transferred $\Ainf$ operations), the MC recursion may not terminate at each arity, and the finiteness argument would need modification.
\item \emph{Non-perturbative completion.}  The resurgent structure (Conjecture~\ref{conj:thqg-I-non-perturbative}) remains open.  The identification of the non-perturbative saddle points (BTZ black holes) is conjectural, and a derivation from first principles (the MC equation) would require extending the shadow theory beyond perturbation theory.
\item \emph{Higher-arity convergence radius.}  The convergence radius $R_{\mathrm{grav}}$ for the full (non-scalar) partition function may exceed $4\pi^2$.  Determining the exact radius requires controlling the growth of higher-arity graph amplitudes, which depends on the shadow depth class.
\item \emph{Analytic sewing beyond HS.}  The HS-sewing condition is sufficient but may not be necessary.  Weaker conditions (e.g., trace-class sewing without the HS factorization) may extend finiteness to a larger class of algebras, including completed W-algebras of infinite rank.
\end{enumerate}
\end{remark}
